\documentclass[a4paper,11pt]{ltjsarticle}

% ========== package ========== %
% --- 数式環境 --- %
\usepackage{amsmath}
\usepackage{amssymb}
\usepackage{mathtools}
% --- 数式環境 END --- %
% ブラケット
\usepackage{braket}
% ベクトルの太文字
\usepackage{bm}
% 添え字左右上下 \fourIdx{左上}{左下}{右上}{右下}{ベース} で使う．Iは大文字
\usepackage{fouridx}
% 色
\usepackage{color}
% 余白指定
\usepackage{geometry}
% --- ページ内リンク --- %
\usepackage{hyperref}
\hypersetup{colorlinks=true,linkcolor=blue,citecolor=blue,urlcolor=blue}
% --- ページ内リンク END --- %
% 花文字
\usepackage{mathrsfs}
% 画像
\usepackage{graphicx}
% tcolorbox
\usepackage{tcolorbox}
% appendix
\usepackage{appendix}
% --- 参考文献 --- %
\usepackage{etoolbox}
\usepackage{fontspec} % fontspecは読み込んでOKですが、設定順序に注意が必要です
%% --- 参考文献フォント --- %%
  \IfFontExistsTF{Times New Roman}{
    \newfontfamily\bibfontface{Times New Roman}
  }{
    \newfontfamily\bibfontface{FreeSerif} % 代替フォント
  }
%% --- 参考文献フォント END --- %%
%% --- 参考文献リンク位置 --- %%
  \BeforeBeginEnvironment{thebibliography}{%
    \phantomsection
    \addcontentsline{toc}{section}{参考文献}
  }
  \AtBeginEnvironment{thebibliography}{%
    \normalfont\bibfontface
  }
%% --- 参考文献のリンク位置 --- %%
% --- 参考文献 END --- %
% ========== package END ========== %

% ========== newcommand ========== %
% --- セクションから目次へのリンク --- %
\newcommand{\sectionlink}[1]{\section[#1]{\hyperlink{toc}{#1}}}
\newcommand{\subsectionlink}[1]{\subsection[#1]{\hyperlink{toc}{#1}}}
% セクション番号リンク
\makeatletter
\renewcommand{\@seccntformat}[1]{%
  \protect\hyperlink{toc}{\csname the#1\endcsname}\quad
}
\makeatother

% 原文の通し番号(1)--(84)を維持するため，節ごとの式番号リセットは行わない．
% ==========- newcommand END ========== %
\allowdisplaybreaks
%===========================================================
% --- タイトル --- %
%===========================================================
\geometry{left=25mm, right=25mm, top=30mm, bottom=30mm}
\def\be{\begin{equation}}
%\def\ee{\end{equation}}
% [[[[[
\def\be{\begin{eqnarray}}
\def\ee{\end{eqnarray}}
\renewcommand{\L}{\mathcal{L}}
\def\d{\partial}
\def\nn{\nonumber}
\newcommand{\dd}[2]{\frac{\partial #1}{ \partial #2}}
% ]]]]]
\def\cK{{\cal{K}}}
\def\cI{{\cal{I}}}
\def\vp{\varphi}
\newcommand{\GeV}{{\rm GeV}}
\newcommand{\TeV}{{\rm TeV}}
\newcommand{\eV}{{\rm eV}}
\newcommand{\half}{{\textstyle{\frac12}}}
\newcommand{\di}{\partial}
\def\ttr{t_{\rm Tr}}
\newcommand{\MP}{M_{\rm Pl}}
\newcommand{\Mbar}{{\,\overline{\vphantom{M}\,\,\,\,\,}\!\!\!\!\!\!M}}



% 原文の数式と長い本文を組むための最小限の追加設定
\setlength{\emergencystretch}{3em}
\renewcommand{\theequation}{\arabic{equation}}
\renewcommand{\theHequation}{\arabic{equation}}
\newcommand{\subsubsectionlink}[1]{\subsubsection[#1]{\hyperlink{toc}{#1}}}
\AtBeginEnvironment{thebibliography}{\interlinepenalty=10000}
\hypersetup{unicode=true,pdftitle={インフレーションの有効場理論},pdfauthor={Clifford Cheung et al.}}
\title{\bfseries インフレーションの有効場理論\\[0.4em]
{\large The Effective Field Theory of Inflation}}
\author{Clifford Cheung$^{\rm a}$，Paolo Creminelli$^{\rm b}$，\\
A.~Liam Fitzpatrick$^{\rm a}$，Jared Kaplan$^{\rm a}$，\\
Leonardo Senatore$^{\rm a}$\\[0.7em]
{\small $^{\rm a}$ Jefferson Physical Laboratory, Harvard University}\\
{\small Cambridge, MA 02138, USA}\\[0.3em]
{\small $^{\rm b}$ Abdus Salam International Center for Theoretical Physics}\\
{\small Strada Costiera 11, 34014 Trieste, Italy}}
\date{\small \href{https://arxiv.org/pdf/0709.0293v2}{arXiv:0709.0293v2} [hep-th]（2008年3月5日）\\IC/2007/032}
\begin{document}
\maketitle
\nocite{Creminelli:2006xe,Alishahiha:2004eh,Chen:2005ad,Shandera:2006ax,Kecskemeti:2006cg,Shiu:2006kj,Arkani-Hamed:2003juy,Dubovsky:2005xd,Cornwall:1974km,Salopek:1990jq,Lyth:2004gb,Cheung:2007sv,Maldacena:2002vr,Seery:2006vu,tHooft:1974toh,Barbieri:2004qk,Hsu:2004vr,Adams:2006sv,Mukhanov:2005bu,Garriga:1999vw,Chen:2006nt,Babich:2004gb,Creminelli:2006rz,Chen:2006dfn,Arkani-Hamed:2003pdi,Arkani-Hamed:2005teg,Senatore:2004rj,Polchinski:1992ed,Buchbinder:2007ad,Creminelli:2007aq,Arkani-Hamed:2007ryv,Wald:1984rg}

\begin{abstract}
単一場模型の場合について，インフレーションの有効場理論，すなわち準de Sitter背景のまわりの揺らぎを記述する最も一般的な理論を研究する．ユニタリーゲージへ移ることにより，スカラー・モードを計量に吸収させることができる．このゲージでは，最も一般的な理論は，$g^{00}$や時間一定面の外的曲率$K_{\mu\nu}$など，空間微分同相変換のもとで不変な最低次元の演算子によって構成される．このアプローチにより，単純なスローロール・インフレーションに対する，ありうるすべての高エネルギー補正を特徴付けることができ，その大きさは実験によって制限される．また，音速の小さい模型やゴースト・インフレーションを含む，すべての単一場模型を共通の言葉で記述し，準de Sitter背景を持つことが何を意味するかを明らかにする．時間微分同相変換の非線形実現によって，音速の低下と非Gaussian性の増大のような，異なる観測量どうしの相関が強制される．
\end{abstract}
\clearpage
\hypertarget{toc}{}
\tableofcontents
\clearpage

\sectionlink{序論}
有効場理論のアプローチ，すなわち基礎となる対称性と両立する最低次元の演算子によって系を記述する方法は，素粒子物理学から物性物理学まで，多くの分野で大きな成果を上げてきた．本論文の目的は，この方法を適用し，インフレーションを起こしている宇宙論的背景のまわりの揺らぎの理論を記述することである．

単一場インフレーション模型を調べる通常の方法では，スカラー場$\phi$のラグランジアンから出発し，$\phi$の運動方程式とFRW計量に対するフリードマン方程式を連立して解く．ここで関心があるのは，スカラー場が一様で時間依存する解$\phi_0(t)$に従い，ハッブル・パラメータがゆっくり変化する加速膨張，すなわちインフレーション解である．そのうえで，この背景解のまわりの摂動を調べ，さまざまな宇宙論的観測量に対する予言を求める．

時間発展する解のまわりの摂動の理論は，出発点とした$\phi$の理論とは大きく異なる．$\phi$はすべての微分同相変換（diffs）のもとでスカラーであるのに対し，摂動$\delta\phi$は空間微分同相変換のもとでのみスカラーであり，時間微分同相変換に対しては非線形に変換する：
\be
t \to t + \xi^0(t, \vec x)   \qquad \delta\phi \to \delta\phi + \dot\phi_0(t) \xi^0 \;.
\ee
とくに，$\phi(t,\vec x)=\phi_0(t)$というゲージを選べば，インフラトンの摂動はなくなり，すべての自由度が計量に含まれる．スカラー変数$\delta\phi$は重力子に吸収され，重力子はスカラー・モードと2つのテンソル・ヘリシティという3つの自由度を持つようになる．この現象は，自発的に破れたゲージ理論で起こることと類似している．ゲージ対称性のもとで非線形に変換するゴールドストーン・モードは，ゲージ・ボソンに吸収され（ユニタリーゲージ），質量を持つスピン1粒子を与える．ゴールドストーン場の非線形シグマ模型は，標準模型のヒッグス・セクターのように，ゲージ対称性の線形表現へ埋め込んで紫外（UV）完備化することができる．これは，微分同相変換を線形に表現する$\phi$のラグランジアンから出発する，通常のインフレーションの定式化に対応する．本論文では，時間微分同相変換が非線形に実現される時間発展する真空のまわりで，インフレーション中の摂動の理論を直接記述するという，もう一つの見方を強調したい．この形式は，ヌル・エネルギー条件を破る可能性を調べるために，文献\cite{Creminelli:2006xe}で一般のFRW背景に対して初めて導入された．ここではインフレーション解に焦点を当てて，この形式を拡張する．

ハッブル率$H(t)$を持つ平坦なFRW背景のまわりの摂動を記述し，時間微分同相変換が破れ，空間微分同相変換は破れていない最も一般的なラグランジアンは，ユニタリーゲージで次のように与えられることを示す：
\begin{eqnarray}
\label{eq:Laguni}
S & = & \int  \! d^4 x  \: \sqrt{- g} \;\Big[\frac12 M_{\rm Pl}^2 R+
M_{\rm Pl}^2 \dot{H} g^{00} -  M_{\rm Pl}^2 \left(3
H^2+\dot{H}\right) + \frac{M_2(t)^4}{2!} (g^{00}+1)^2 \\ & & \nonumber + \frac{M_3(t)^4}{3!} (g^{00}+1)^3 + \ldots
-\frac{\bar{M_2}(t)^2}{2}\delta K^\mu {}_\mu {}^2 +... \Big] \; .
\end{eqnarray}
アインシュタイン・ヒルベルト項の次にある2つの演算子は，指定された非摂動解$H(t)$を持つという要請によって固定される．一方，それ以外の演算子はすべて自由であり，同じ背景解を持つさまざまな摂動の理論をすべてパラメータ化する．時間微分同相変換が破れているので，たとえば$g^{00}$や時間一定面の外的曲率$K^\mu {}_\nu {}$など，空間微分同相変換を尊重する任意の項を書くことができる．演算子の係数は一般に時間に依存する．微分同相変換のもとで不変でないラグランジアンを使うことに，読者は不安を感じるかもしれない．しかし，通常のゲージ理論と同様に，微分同相不変性は明らかに回復できる．パラメータ$\xi^0(t,\vec x)$を持つ時間微分同相変換を行い，そのパラメータを，時間微分同相変換のもとで$\pi(t,\vec{x})\rightarrow \pi(t,\vec{x})-\xi^0(t,\vec{x})$とシフトする場$\pi(t, \vec x)$へ昇格させればよい．スカラー$\pi$は時間微分同相変換を非線形に実現するゴールドストーン・モードであり，FRW解のまわりのスカラー摂動を記述する．

質量を持つゲージ・ボソンの縦成分の物理は，十分高いエネルギーでは，スカラーのゴールドストーン・モードに着目して調べられることがよく知られている（等価定理）．ここでも同じことが成り立つ．十分高いエネルギーでは重力との混合は重要でなくなり，ゴールドストーン・モードに集中できる．この領域では物理が非常に見通しよくなり，宇宙論的摂動についての情報の大部分を得ることができる．ラグランジアン(\ref{eq:Laguni})に対して破れた微分同相変換を行い，ゴールドストーン・モード$\pi$に着目すると，次を得る：
\begin{eqnarray} \label{theeffectiveaction} &&S_\pi = \int \! d^4 x \:
\sqrt{-g}\left[ M_{\rm Pl}^2 \dot H
 \, (\partial_\mu \pi)^2\right.\nonumber \\ &&\qquad\left.+2 M^4_2\left(\dot\pi^2+ \dot\pi^3-\dot\pi \frac{1}{a^2}(\di_i\pi)^2\right)-\frac{ 4}{3} M^4_3 \dot\pi^3 - \frac{\bar M ^2}{2} \,
\frac{1}{a^4}(\di_i ^2 \pi)^2 +\ldots
\right].
 \end{eqnarray}

ユニタリーゲージにおけるそれぞれの不変演算子は，ゴールドストーン場の（非線形な）演算子へと昇格する．微分同相不変性の非線形実現によって，各種の項のあいだに関係が強制されるのである．

理論の系統的な構成に進む前に，このアプローチの利点を簡単に述べておこう．
\begin{itemize}
\item 標準的な運動項とスローロール・ポテンシャルを持つスカラー場による，最も標準的なインフレーションのシナリオから出発し，ありうるすべての高エネルギー効果についてのわれわれの無知を，主要な不変演算子によってパラメータ化した．たとえば摂動の非ガウス性の測定や，重力波のスペクトルの傾きに対する整合性関係からのずれを調べることによって，実験は各演算子に制限を与える．ある意味で，これは素粒子物理学で行うことと類似している．そこでは，標準模型からのずれを記述し，したがって新しい物理の効果を表す演算子の大きさに制限を課す．

\item 対称性と，背景がインフレーションを起こすという要請によって何が強制され，何が自由であるかが明確になる．たとえば式(\ref{theeffectiveaction})は，空間運動項$(\nabla \pi)^2$が$\dot H$に比例する一方，時間運動項$\dot\pi^2$は自由であることを示している．もう一つの例は，ユニタリーゲージの演算子$(g^{00}+1)^2$である．これをゴールドストーン場$\pi$で書くと，$\pi$励起の音速を低下させる2次項$\dot\pi^2$と，モード間の相互作用，すなわち非ガウス性を増大させる3次項$\dot\pi (\nabla\pi)^2$が生じる．したがって，他の演算子との相殺がなければ，音速の低下と非ガウス性の増大は対称性によって結び付けられる．さらに，$\pi$励起が光円錐の外へ伝播することを避けるため，この演算子の係数は正でなければならない．

\item 可能な演算子をすべて把握できる．たとえば微分の最低次では，3つの$\pi$モード間の相互作用は$(g^{00}+1)^2$と$(g^{00}+1)^3$によって変更できる．これは3点関数の異なる2つの形状に対応し，原理的には実験によって区別して，各演算子の大きさを決めることができる．

\item 可能な単一場模型がすべて統一される．たとえば，かなり特徴的な予言を持つDBIインフレーション\cite{Alishahiha:2004eh,Chen:2005ad,Shandera:2006ax,Kecskemeti:2006cg,Shiu:2006kj}のように，ラグランジアンを$L((\partial\phi)^2,\phi)$へ変更した模型が関心を集めてきた．われわれの記述では，これらは演算子$(g^{00}+1)^n$が大きい場合に対応する．もう一つの興味深い極限は$\dot H \to 0$である．この場合，主要な空間運動項は$\bar M^2$に比例する演算子から生じ，$(\nabla^2 \pi)^2$という形になる．この極限はゴースト・インフレーション\cite{Arkani-Hamed:2003uz}を記述する．

\item $\phi$による記述では，場の再定義$\phi \to \tilde\phi(\phi)$を行うことができる．得られるラグランジアンが同じ物理を記述することは確かだが，そのことは自明ではない．簡単な例は，ラグランジアン
\be
f(\phi)^2 (\partial\phi)^2 -V(\phi) \;,
\ee
で与えられる．ここで$f$は一般の関数である．これは，場の再定義$\tilde\phi(\phi)$，$d\tilde\phi/d\phi = f(\phi)$を通じて，標準的な運動項と別のポテンシャルを持つラグランジアンに等価である．しかし，より一般的な項を考えると，異なるラグランジアン間の等価性はさらに複雑になる．一方，時間微分同相変換の一種の標準的な非線形表現を実現する$\pi$のレベルには，この曖昧さは存在しない．

\item $\phi$による記述では，演算子に含まれる場の一部を背景解上で評価できるため，摂動を調べる際に演算子の重要性をどう評価すべきかは明らかでない．たとえば$(\partial\phi)^{2n}$という形のすべての演算子を含む理論では，DBIインフレーションで起こるように，背景の速度$\dot\phi_0$が十分大きければ，それらがすべて同じ重要性を持ちうる．一方，$\pi$による記述では，演算子の重要性を見積もる通常の方法が機能する．このことはゴースト・インフレーションではさらに明瞭であり，$\pi$が非相対論的な分散関係を持つため，演算子のスケーリングは$\pi$による記述ではじめて明確になる．

\item 宇宙論的摂動のループ補正を計算する際には，解のまわりで演算子を直接パラメータ化することが決定的に重要である．外線数が与えられたダイアグラムは，一般にUV発散を含む．ラグランジアン(\ref{eq:Laguni})では，$n$個の摂動間の相互作用を記述する項が有限個しかないので，この発散は容易に繰り込める．一方，$\phi$のラグランジアンのレベルでは，$n$個の摂動間の相互作用に寄与する演算子が\emph{無限個}存在する．実際，各演算子の多数の脚を背景上に置くことができるため，演算子と摂動のダイアグラムとの関係はかなり不明瞭になる．
\end{itemize}

\sectionlink{ユニタリーゲージにおける作用の構成}
インフレーションとは，ハッブル・パラメータがほぼ一定である宇宙の加速膨張期である．この準デ・ジッター背景には，物理的な時計によって定まる特別な空間的スライシングがあり，これによって減速膨張する高温ビッグバンの発展へ滑らかに接続できる．このスライシングは通常，時間発展するスカラー$\phi(t)$によって実現される．念頭に置いておくとよいもう一つの例は完全流体である\footnote{実際，たとえば文献\cite{Dubovsky:2005xd}で示されているように，完全流体の渦なしの励起は，微分結合するスカラーによって記述できる．}．この解のまわりの摂動を記述するため，特別なスライシングが$t$一定面と一致するゲージ，すなわち$\delta\phi(\vec x,t)=0$を選ぶことができる．このゲージでは明示的なスカラー摂動はなく，計量の揺らぎだけが存在する．時間微分同相変換は固定され，もはやゲージ対称性ではないので，重力子は3つの自由度を記述する．スカラー摂動は計量に吸収されたのである．

このゲージにおける最も一般的なラグランジアンは何だろうか．計量$g_{\mu\nu}$の関数であり，（線形に実現された）時間依存する空間微分同相変換$x^i\rightarrow x^i+\xi^i(t,\vec{x})$のもとで不変な演算子を書き下さなければならない．実際，空間微分同相変換は破れていない．すべての微分同相変換のもとで不変なリーマン・テンソルを含む通常の項に加え，系の対称性が小さくなっているため，さらに多くの項が許される．それらは，重力子に吸収された追加の自由度を記述する．たとえば，$g^{00}$が空間微分同相変換のもとでスカラーであり，ユニタリーゲージのラグランジアンに自由に現れうることは容易に分かる．微分を含まない項は，$g^{00}$の多項式に限られる．時空に特別なスライシングがあるので，そのスライシングを記述する幾何学的な量を書くことも許される．たとえば，時間一定面の外的曲率$K_{\mu\nu}$は空間微分同相変換のもとでテンソルであり，作用中に用いることができる．作用のどの項にも，一般の時間の関数を掛けられることに注意しよう．\ref{app:generic}では，最も一般的なラグランジアンが次のように書けることを証明する：
\begin{eqnarray}
\label{eq:action}\nonumber
S & = & \int  \! d^4 x \; \sqrt{- g} \Big[ \frac12 M_{\rm Pl}^2 R - c(t)
g^{00} -  \Lambda(t) + \frac{1}{2!}M_2(t)^4(g^{00}+1)^2+\frac{1}{3!}M_3(t)^4 (g^{00}+1)^3+ \\
&& - \frac{\bar M_1(t)^3}{2} (g^{00}+1)\delta K^\mu {}_\mu
-\frac{\bar M_2(t)^2}{2} \delta K^\mu {}_\mu {}^2
-\frac{\bar M_3(t)^2}{2} \delta K^\mu {}_\nu \delta K^\nu {}_\mu + ...
\Big] \; ,
\end{eqnarray}
ここで省略記号は，揺らぎについて高次の項，またはより多くの微分を持つ項を表す．非摂動FRW背景からの時間一定面の外的曲率の変化を$\delta K_{\mu\nu}$と書く．すなわち，$\delta K_{\mu\nu} = K_{\mu\nu} - a^2 H h_{\mu\nu}$であり，$h_{\mu\nu}$は誘導空間計量である．上の作用で，選んだFRW解のまわりの1次の摂動を含むのは最初の3項だけであり，他の項はすべて明示的に2次以上であることに注意しよう．したがって，係数$c(t)$と$\Lambda(t)$は，指定されたFRWの発展$H(t)$を持つという要請，すなわちこの解のまわりでタドポール項が相殺するという要請によって固定される．これらの係数を固定する前に，この単純化が自明ではないことを理解しておくことが重要である．背景解のまわりで1次の寄与を与える演算子は，無限個あると予想されるかもしれない．しかし作用は，$\delta K_{\mu\nu}$や$g^{00}+1$のような1次の項の多項式として書けるので，演算子が1次，2次，あるいはそれ以上の次数から始まるかが明確になる．式(\ref{eq:action})の項以外のすべての1次項は微分を含み，部分積分によって，ここで考えた3つの1次項と高次の共変な項との組合せにできる．この構成は\ref{app:tadpoles}で具体的に行う．したがって，\emph{非摂動の膨張史によって$c(t)$と$\Lambda(t)$が固定され，異なる模型間の違いは高次の項に表される}．

指定されたFRWの発展が解であることを課して，1次項を固定しよう．すでに述べたように，$c$と$\Lambda$に比例する項だけが，摂動の0次で消えず，したがってアインシュタイン方程式の右辺に寄与するエネルギー運動量テンソル
\begin{equation}
T_{\mu \nu} = -\frac{2}{\sqrt{-g}}\frac{\delta S_{\rm
matter}}{\delta g^{\mu\nu}}
\end{equation}
を与える．インフレーション中に主として関心があるのは，平坦なFRW宇宙である（一般の場合については\ref{app:tadpoles}を参照）：
\be
ds^2 = -dt^2 + a^2(t) d \vec{x}^2 
\ee
このため，フリードマン方程式は次で与えられる：
\begin{eqnarray}
H^2 & = & \frac{1}{3 M_{\rm Pl}^2} \big[ c(t)+\Lambda(t)\big]  \\
\frac{\ddot a}{a} = \dot H + H^2 & =  & -\frac{1}{3 M_{\rm Pl}^2} \big[ 2
c(t)-\Lambda(t) \big] \;.
\end{eqnarray}
$c$と$\Lambda$について解くと，作用(\ref{eq:action})を次のように書き換えられる：
\begin{eqnarray}
\label{eq:actiontad}
S &=& \int  \! d^4 x \; \sqrt{- g} \Big[ \frac12 M_{\rm Pl}^2 R + M_{\rm Pl}^2 \dot H
g^{00} - M_{\rm Pl}^2 (3 H^2 + \dot H) +  \nonumber \\ && \frac{1}{2!}M_2(t)^4(g^{00}+1)^2+\frac{1}{3!}M_3(t)^4 (g^{00}+1)^3+ \nonumber \\
&& - \frac{\bar M_1(t)^3}{2} (g^{00}+1)\delta K^\mu {}_\mu
-\frac{\bar M_2(t)^2}{2} \delta K^\mu {}_\mu {}^2
-\frac{\bar M_3(t)^2}{2} \delta K^\mu {}_\nu \delta K^\nu {}_\mu + ...
\Big] \; .
\end{eqnarray}
すでに述べたように，上の作用に含まれる演算子の係数はすべて，一般の時間依存性を持ちうる．しかし，ここで関心があるのは，1ハッブル時間のあいだに$H$と$\dot H$が大きく変化しない解である．したがって，他のすべての演算子についても同じことが成り立つと仮定するのが自然である．この仮定のもとで，ラグランジアンは近似的に時間並進不変になる\footnote{時間シフトが厳密な対称性となる極限では$\dot H \to 0$となるため，注意して極限を取らなければならない．ゴースト・インフレーションの議論で見るように，これはゴールドストーン場の空間運動項が消えることを意味する．}．したがって，ループ効果によって生じる時間依存性は，小さな対称性の破れのパラメータによって抑制される\footnote{この対称性は，時間微分同相変換の破れとは無関係であることに注意しよう．$\phi$による記述でこの対称性がどう現れるかを見るには，適切な場の再定義を行えば，常に$\dot\phi =$ const.と仮定できることに注目すればよい．この選択のもとでは，シフト対称性$\phi \to \phi $ + const.から，ユニタリーゲージのラグランジアンの時間並進不変性が従う．この対称性と$\phi$のラグランジアンの時間並進対称性は，背景によって対角部分群へと破られる．この残存対称性が，ユニタリーゲージのラグランジアンにおける時間シフトである．}．この仮定はとくに都合がよい．係数が急速に時間変化すると，指数関数的膨張が生む摩擦に打ち勝つことがあり，その結果，標準的なFRW宇宙論の一様性問題を解決するために必要な力学的アトラクターとしての性質を，インフレーションが失うおそれがあるからである．

このアプローチが記述するのは，スカラー・モードだけでなく重力についても最も一般的なラグランジアンであることを強調しておく．たとえば，高エネルギー効果はリーマン・テンソルの摂動$\delta R_{\mu\nu\rho\sigma}$を含む演算子に表される．これらの補正は微分について高次であるため，以下では明示的には扱わない．

単純なインフレーション模型をこの記述でどのように書くか，いくつか例を挙げよう．標準的な運動項とスローロール・ポテンシャル$V(\phi)$を持つ模型は，ユニタリーゲージで次のように書ける：
\be
\int \! d^4x \: \sqrt{-g}
\left[ -\frac 1 2 (\partial \phi)^2 - V(\phi) \right]  \to  \int \!
d^4x \: \sqrt{- g} \left[ -\frac{\dot \phi_0(t)^2}{2} g^{00} - V(\phi_0(t))
\right] \; .
\ee
フリードマン方程式から$\dot\phi_0(t)^2=-2M^2_P \dot{H}$および$V(\phi(t))=M_{\rm Pl}^2 (3H^2+\dot H$)が得られるので，作用は，最初の3項以外をすべてゼロとした式(\ref{eq:actiontad})の形であることが分かる．ただし，他のすべての項はループ補正によって生成されるため，これが厳密に成り立つことは明らかにない．それらの項は，この単純なスローロール・インフレーション模型に対する高エネルギー物理の，ありうるすべての効果を表している．

さらに一般的な場合として，各$\phi$に作用する微分が高々1つであるすべてのラグランジアンを含めることができる．すなわち，$L= P(X,\phi)$，ただし$X=g^{\mu\nu}\partial_\mu\phi\partial_\nu\phi$である．非摂動解$\phi_0(t)$のまわりでは，
\be
S = \int \! d^4x \: \sqrt{-g} \;  P(\dot\phi_0(t)^2 g^{00}, \phi(t))
\ee
となる．これは明らかに先ほどの形であり，$M_n^4(t) =\dot\phi_0(t)^{2n} \partial^n P/\partial X^n$を$\phi_0(t)$で評価すればよい．外的曲率を含む項では，1つのスカラーに2つ以上の微分が作用し，これらの項は厳密なデ・ジッター極限$\dot H \to 0$で決定的に重要になる．


\sectionlink{ゴールドストーン・ボソンの作用\label{sec:Goldstone}}
ユニタリーゲージのラグランジアンは，重力子の2つのヘリシティと1つのスカラー・モードという3つの自由度を記述する．破れた時間微分同相変換を行うと（シュテュッケルベルクの方法），このスカラー・モードは，その対称性を非線形に実現するゴールドストーン・ボソンとして明示的に現れる．質量を持つゲージ・ボソンの縦成分に対する等価定理\cite{Cornwall:1974km}との類推から，混合を無視できる短距離では，ゴールドストーン場の物理は重力子の2つのヘリシティから分離すると予想される．同じ方法をここで適用する前に，非アーベル・ゲージ理論で何が起こるかを簡単に振り返ろう．

非アーベル・ゲージ場$A_\mu^a$のユニタリーゲージにおける作用は
\be
S = \int \! d^4x  \,-\frac{1}{4} {\rm Tr}\, F_{\mu\nu}F^{\mu\nu}-\frac12 m^2 {\rm Tr}\, A_\mu A^\mu \ ,
\ee
であり，ここで$A_\mu = A_\mu^a T^a$である．ゲージ変換のもとで
\be
\label{eq:AmuU}
A_\mu \to U A_\mu U^\dagger + \frac{i}{g} U \partial_\mu U^\dagger \equiv \frac{i}{g} U D_\mu U^\dagger \;.
\ee
となる．したがって作用は
\be
S = \int \! d^4x  \,-\frac{1}{4} {\rm Tr}\, F_{\mu\nu}F^{\mu\nu} - \frac12 \frac{m^2}{g^2} {\rm Tr} D_\mu U^\dagger D_\mu U \;.
\ee
となる．$U=\exp{[i T^a \pi^a(t,\vec x)]}$と書くことによって，ゲージ不変性を「回復」できる．ここで$\pi^a$はスカラー（ゴールドストーン場）であり，ゲージ変換$\Lambda$のもとで次のように非線形に変換する：
\be
e^{i T^a \widetilde\pi^a(t, \vec x)} = \Lambda(t, \vec x) \,e^{i T^a \pi^a(t,\vec x)}
\ee
正準規格化$\pi_c \equiv m/g \cdot \pi$へ移ると，ゴールドストーン・ボソンの自己相互作用はスケール$4 \pi m/g$で強結合になることが分かる．このスケールは，パラメータ依存性の意味でゲージ・ボソンの質量より高い．ゴールドストーン場を再導入する利点は，エネルギー$E \gg m$で，ゴールドストーン場とゲージ場の横成分との混合が重要でなくなり，2つのセクターが分離することである．実際，式(\ref{eq:AmuU})の混合項は
\be
\frac{m^2}{g} A_{\mu}^a \partial^\mu \pi^a = m A_{\mu}^a \partial^\mu \pi_c^a
\ee
という形を取り，$E \gg m$では正準運動項$(\partial \pi_c)^2$に比べて重要でなくなる．領域$m \ll E \ll 4 \pi m /g$では，ゴールドストーン場$\pi$の物理は弱結合であり，横成分との混合を無視して調べることができる．

時間微分同相変換が破れた場合について，同じ手順をたどろう．たとえば，次の2つの演算子に着目する：
\begin{equation}
\int d^4x\;  \sqrt{-g} \left[A(t)+B(t)g^{00}(x)\right] \ .
\end{equation}
破れた時間微分同相変換$t \to \widetilde t= t + \xi^0(x)$，$\vec{x} \to \vec{\widetilde{x}}=\vec{x}$のもとで，$g^{00}$は次のように変換する：
\begin{equation}
g^{00}(x)\to \widetilde g^{00}(\widetilde x(x))=\frac{\partial \widetilde x^0(x)}{\partial x^\mu}\frac{\partial \widetilde x^0(x)}{\partial x^\nu} g^{\mu\nu}(x) \,  .
\end{equation}
変換後の場で書いた作用は次で与えられる：
\begin{eqnarray}
\int d^4x\;  \sqrt{-\widetilde g(\widetilde x(x))} \left|\frac{\partial \widetilde x}{\partial x} \right| \left[A(t)+B(t) \frac{\partial x^0}{\partial \widetilde x^\mu}\frac{\partial x^0}{\partial \widetilde x^\nu} \widetilde g^{\mu\nu}(\widetilde x(x))\right]\ .
\end{eqnarray}
積分変数を$\widetilde x$へ変更すると，
\begin{eqnarray}
&&\int d^4\widetilde x\;  \sqrt{-\widetilde g(\widetilde x)} \left[A(\widetilde t-\xi^0(x(\widetilde x)))\right.\nonumber \\ &&\qquad\left.+B(\widetilde t-\xi^0(x(\widetilde x))) \frac{\partial (\widetilde t-\xi^0(x(\widetilde x)))}{\partial \widetilde x^\mu}\frac{\partial (\widetilde t-\xi^0(x(\widetilde x)))}{\partial \widetilde x^\nu} \widetilde g^{\mu\nu}(\widetilde x)\right].
\end{eqnarray}
を得る．ゴールドストーン場を再導入する手続きは，ゲージ理論の場合と同様である．上の作用に現れる$\xi^0$をすべて，
\begin{equation} 
\xi^0(x(\widetilde x)) \to - \widetilde \pi(\widetilde x ) \, .
\end{equation}
と置き換える．簡単のためチルダを省略すると，次を得る：
\begin{eqnarray}
\int d^4x\;  \sqrt{- g(x)} \left[A( t+\pi(x))+B(t+\pi(x)) \frac{\partial (t+\pi(x))}{\partial x^\mu}\frac{\partial (t+\pi(x))}{\partial x^\nu} g^{\mu\nu}(x)\right].
\end{eqnarray}
$\pi$に変換則
\begin{equation}
\pi(x) \to \widetilde\pi(\widetilde x(x))=\pi(x)-\xi^0(x) \ .
\end{equation}
を割り当てれば，上の作用が無限小変換だけでなく，すべての次数で微分同相変換のもとで不変となることを確認できる．この定義のもとで，$\pi$はスカラー場として変換するとともに，時間微分同相変換に対して追加のシフトを受ける．

この手続きをユニタリーゲージの作用(\ref{eq:actiontad})に適用すると，
\begin{eqnarray}\label{Smixed}
S = \int \! d^4 x  \: \sqrt{- g} &&\left[\frac{1}{2}M_{\rm Pl}^2 R
- M^2_{\rm Pl} \left(3H^2(t+\pi) +\dot{H}(t+\pi)\right)+ \right.\\ \nn
&&+M^2_{\rm
Pl} \dot{H}(t+\pi)\left(
(1+\dot\pi)^2g^{00}+2(1+\dot\pi)\partial_i\pi g^{0i}+
g^{ij}\partial_i\pi\partial_j\pi\right) + \\ \nonumber
&&\frac{M_2(t+\pi)^4}{2!}\left(
(1+\dot\pi)^2g^{00}+2(1+\dot\pi)\partial_i\pi g^{0i}+
g^{ij}\partial_i\pi\partial_j\pi+1\right)^2 + \nonumber\\
\nonumber && \left. \frac{M_3(t+\pi)^4}{3!}\left(
(1+\dot\pi)^2g^{00}+2(1+\dot\pi)\partial_i\pi g^{0i}+
g^{ij}\partial_i\pi\partial_j\pi+1\right)^3+ ... \right] \; ,
\end{eqnarray}
を得る．ここではひとまず，簡単のため外的曲率を含む項を無視した．

この作用はかなり複雑であり，この時点では，ユニタリーゲージのラグランジアンからゴールドストーン場$\pi$を再導入する利点は明らかでない．ゲージ理論の場合と同様に，十分短い距離では計量の揺らぎを無視してゴールドストーン場の物理を調べられるため，単純化が起こる．それが可能な領域は，ゲージ理論の場合と同じく，上のラグランジアンの混合項を見るだけで見積もれる．実際，式(\ref{Smixed})では，$\pi$と$g_{\mu\nu}$を混合する2次項の微分の数は$\pi$の運動項より少なく，ある高エネルギー・スケールより上で無視できることが分かる．一般に，その答えはどの演算子が存在するかに依存する．まず，タドポール項だけが重要な最も簡単な場合（$M_2=M_3=\ldots=0$）から始めよう．これは通常のスローロール・インフレーションに対応する．重力との主要な混合は，次の形の項から生じる：
\begin{equation}
\sim M_{\rm Pl}^2 \dot H \dot\pi \delta g^{00} \ .
\end{equation}

正準規格化（$\pi_c\sim M_{\rm Pl} \dot H^{1/2}\pi,\  \delta g_c^{00}\sim M_{\rm Pl} \delta g^{00}$）を行うと，$E_{\rm mix}\sim \epsilon^{1/2} H$より高いエネルギーでは混合項を無視できることが分かる．ここで$\epsilon$は通常のスローロール・パラメータ$\epsilon\equiv -\dot H/H^2$である．もう一つ関心があるのは，$M_2$の演算子が大きくなる場合である．この場合には，
\begin{equation}
\sim M_2^4 \dot \pi \delta g^{00}
\end{equation}
という形の混合項があり，正準規格化を行うと（今度は$\pi_c\sim M_2^2\pi$であることに注意），$E_{\rm mix}\sim M_2^2/M_{\rm Pl}$より高いエネルギーで無視できるようになる\footnote{ここで調べている理論ではローレンツ対称性が自発的に破れているので，混合を無視できるエネルギーと運動量の領域をそれぞれ別に定義すべきである．宇宙論的摂動については，エネルギーの範囲だけに関心を持つ．}．

領域$E\gg E_{\rm mix}$では，作用は大幅に単純化され，次のようになる：
\begin{eqnarray}\label{Spi} 
&&\! \! S_{\rm \pi} = \int \!  d^4 x   \sqrt{- g} \left[\frac12 M_{\rm Pl}^2 R -M^2_{\rm Pl}
\dot{H} \left(\dot\pi^2-\frac{ (\partial_i \pi)^2}{a^2}\right)
\right.\nonumber \\ &&\qquad\left.+2 M^4_2
\left(\dot\pi^2+\dot{\pi}^3-\dot\pi\frac{(\partial_i\pi)^2}{a^2}
\right) -\frac{4}{3} M^4_3 \dot{\pi}^3+ ... \right].
\end{eqnarray}

インフレーション模型が与えられたとき，現在の宇宙論的観測に対する予言を計算したい．この観点からは，デカップリング極限(\ref{Spi})は，このような極端な赤外スケールにはまったく無関係に思える．しかし，通常の単一場スローロール・インフレーションの場合と同じく，摂動論の任意の次数で地平線の外側において一定となる量，通常の$\zeta$変数が存在することを証明できる\cite{Salopek:1990jq,Lyth:2004gb}（高階の空間微分を持つ項を含めた一般化については，文献\cite{verification}の付録Dを参照）．保存量が存在する直観的な理由は，地平線を出た後には，異なる領域がまったく同じように発展することにある．唯一の違いは，ある領域が別の領域に比べてどれだけ膨張したかという点であり，この差が一定に保たれる．

したがって，問題は地平線通過直後の相関関数を計算することに帰着する．そこで，$H$程度の赤外（IR）エネルギー・カットオフを持つラグランジアンを調べればよい．デカップリング・スケール$E_{\rm mix}$が$H$より小さければ，$\pi$のラグランジアン(\ref{Spi})は，$E_{\rm mix}/H$で抑制された項を除いて正しい予言を与える．

すでに述べたように，ユニタリーゲージのラグランジアンの係数の時間依存性は，ハッブル時間に比べて緩やかであり，何らかの一般化されたスローロール・パラメータによって抑制されると仮定している．これは，係数のテイラー展開から生じる追加の$\pi$の項が小さいことを意味する．とくに，3次項のような関連演算子，すなわち赤外へ向かうにつれて支配的になる演算子は，スケール$H$では重要でなく，そのためラグランジアン(\ref{Spi})では無視してある．

以上から，ラグランジアン(\ref{Spi})を用いれば，重力との混合に支配されないすべての観測量を計算できる．混合に支配される量の例には，通常のスローロール・インフレーションにおける非ガウス性がある\cite{Maldacena:2002vr,Seery:2006vu}．ただし，スペクトルの傾きは，主要次数ではラグランジアン(\ref{Spi})から計算できることに注意しよう．後で見るように，その値は実際，混合項を無視して計算した地平線通過時のパワー・スペクトルから単純に導ける．重力との混合が重要になり，ゴールドストーン場の作用だけでは予言に不十分となる場合にも，われわれのアプローチが妥当性を失うわけではないことを強調しておく．作用(\ref{eq:actiontad})は模型についてのすべての情報を含んでおり，重力との混合が大きい場合にも，これを用いてすべての予言を計算できる．


\sectionlink{単一場インフレーションのさまざまな極限}

\subsectionlink{スローロール・インフレーションと高エネルギー補正}
一般的なラグランジアン (\ref{eq:actiontad}) の最も単純な例は，背景のHubbleパラメータ $H(t)$ がわかれば決まる最初の三つの項だけを残し，それ以外の高次の演算子をすべてゼロ，すなわち $M_2 = M_3 = \bar M_1 =\bar M_2 \ldots =0$ とすることで得られる．$\phi$ を用いた記述では，これは高次の項を含まない標準的なスローロール・インフレーションに対応する．
この場合，前節で論じたように，重力との混合を無視してGoldstoneモードのラグランジアン (\ref{Spi}) に着目することで，スケール $H$ における予言を行える．計算したいのは，地平線通過の直後における保存量 $\zeta$ である．これは線形の次数では，ゲージを $\pi = 0$ （本論文の用語ではユニタリーゲージ）に選び，計量の空間成分を
\be
\label{eq:zetagauge}
g_{ij} = a^2(t) \left[(1+ 2 \zeta(t,\vec x)) \delta_{ij} + \gamma_{ij}\right]
\ee
とすることで定義される．ここで $\gamma$ は横波かつトレースレスであり，二つのテンソル自由度を記述する．
$\pi$ と $\zeta$ の関係は非常に単純である．重力との混合を無視しているので，$\pi$ を用いた記述では計量に摂動はない．$\pi = 0$ とするには時間微分同相変換 $t \to t - \pi(t,\vec x)$ を行えばよく，これによって (\ref{eq:zetagauge}) の形の空間計量が得られ，
\be
\label{eq:pizeta}
\zeta (t,\vec x) = - H \pi(t, \vec x) \;.
\ee
となる．

各モード $k$ について，関心があるのは地平線通過 $\omega(k) = k/a \sim H$ の前後の力学だけである．この期間には，スローロール補正を除いて背景をde Sitter時空で近似できる．したがって，正準規格化されたスカラー $\pi_c$ の2点関数は，de Sitter時空での結果
\be
\langle \pi_c(\vec k_1) \pi_c(\vec k_2)\rangle = (2 \pi)^3 \delta(\vec k_1 + \vec k_2) \frac{H^2_*}{2 k_1^3} \;,
\ee
で与えられる．ここでも以下でも，$*$ は地平線通過時の量の値を表す．これから，$\zeta$ の2点関数は
\be
\label{eq:spectrumzeta}
\langle \zeta(\vec k_1) \zeta(\vec k_2)\rangle = (2 \pi)^3 \delta(\vec k_1 + \vec k_2) \frac{H^4_*}{4 M_{\rm Pl}^2 |\dot H_*|}\frac{1}{k_1^3} = (2 \pi)^3 \delta(\vec k_1 + \vec k_2) \frac{H_*^2}{4 \epsilon_* M_{\rm Pl}^2} \frac{1}{k_1^3} \;.
\ee
で与えられることがわかる．変数 $\zeta$ は地平線の外側で一定なので，この式はスローロール補正を除いてすべての $k$ に対して厳密である．特に，これを用いればスローロールの最低次でスペクトルの傾きを計算でき，
\be
n_s-1= \frac{d}{d \log k} \log \frac{H_*^4}{|\dot H_*|} = \frac1{H_*} \frac{d}{d t_*} \log \frac{H_*^4}{|\dot H_*|} = 4 \frac{\dot H_*}{H_*^2} - \frac{\ddot H_*}{H_* \dot H_*} \;.
\ee
となる．

ただし，すべての観測量が $\pi$ のラグランジアン (\ref{Spi}) から計算できるわけではないことに注意しよう．最低次の結果が重力との混合に由来する場合，あるいはスローロール展開の高次で生じる場合には，この方法では計算できない．例えば，式~(\ref{Spi}) の最初の二つの項は $\pi$ の自己相互作用を含まないため，3点関数 $\langle \zeta(\vec k_1) \zeta(\vec k_2) \zeta(\vec k_3) \rangle $ はゼロになってしまう．したがって，式~(\ref{Smixed}) における重力との混合を考慮して，次の次数の補正を調べる必要がある．

高次の項をすべてゼロにするという選択が，厳密には成り立ち得ないことは明らかである．ツリーレベルでゼロにしたとしても，少なくとも輻射補正によってそれらは生成される．この理論は繰り込み不可能であり，摂動展開を十分高次まで進めると，あらゆる相互作用が発散する係数を伴って生成される．追加の項は重力子のループによって生成されるので，非常に小さいこともあり得る．例えば，ユニタリーゲージの相互作用 $M_{\rm Pl}^2 \dot H g^{00}$ から出発すると，$(g^{00}+1)^2$ の形の項が，対数的に発散する係数 $M_2^4 \sim \dot H^2 \log
\Lambda$ を伴って生成されることは容易に確かめられる．これは，$M^4_2 \gtrsim \dot H^2$ と仮定すべきことを意味する（\footnote{重力に結合した質量ゼロのスカラー場の理論における対数発散の具体的な計算は，かなり以前に文献 \cite{tHooft:1974toh} で行われている．}）．ただし，この下限は非常に小さい．例えば，$\pi$ の分散関係は，時間方向の運動項への追加の寄与によって変化する．以下で詳しく論じるように，これは $\pi$ の励起の速度が光速からわずかにずれ，その相対的なずれが $1-c_s \sim M_2^4/(|\dot H| M_{\rm Pl}^2) \sim |\dot H|/M_{\rm Pl}^2$ となることを意味する．スカラーのスペクトルの規格化，式~(\ref{eq:spectrumzeta}) を用いると，光速からのずれは $\gtrsim \epsilon^2 \cdot 10^{-10}$ であることがわかる．あまり興味深い下限ではない．

Planckスケールよりも低いスケールで新たな物理が加わるなら，追加の演算子の大きさはずっと大きくなる．一般に，本論文の手法は，新しい物理がもたらし得るあらゆる効果を正しくパラメータ化する．有効場理論の手法では通常そうであるように，模型の紫外完備化の詳細は高次元演算子に符号化される．これは，標準模型を超える物理で起こることと非常によく似ている．低エネルギーでは，新しい物理がもたらし得る効果は，対称性と両立する一連の高次元演算子に符号化される \cite{Barbieri:2004qk}．標準模型を実験で詳しく調べることによって，これらの高次元演算子の大きさに厳しい制限を課すことができる．本論文の場合にも同じことができるが，残念ながら，考えられる観測の種類ははるかに限られる．
高次元演算子に対して課し得る実験的制限の一例は，重力波のスペクトルの傾きに関する整合性関係である．よく知られているように，Einstein--Hilbert作用から得られる重力波のスペクトルは，
\be
\label{eq:GWspectrum}
\langle \gamma^s(\vec k_1) \gamma^{s'}(\vec k_2)\rangle = (2 \pi)^3
\delta(\vec k_1 + \vec k_2) \frac{H_*^2}{M_{\rm Pl}^2}\frac{1}{k_1^3}
\delta_{s s'}\,
\ee
で与えられる．ここで $\gamma^s$ は，重力波の二つの偏極を表す．この寄与とスカラーの寄与 (\ref{eq:spectrumzeta}) の比は $\epsilon_*$ で与えられる．したがって，重力波のスペクトルの傾き $n_g=-2 \epsilon_*$ は，テンソルモードとスカラーモードの比がわかれば決まる．

この予言は，$M_2 =0$，{\em すなわち} $c_s=1$ を仮定する場合に有効である．実際，後で見るように，スカラーのスペクトルは $c_s^{-1}$ に比例する一方，重力波に関する予言は $M_2$ によって変わらない．整合性関係が実験的に検証されれば，たとえ誤差が大きくても，$c_s$ が $1$ から大きくずれ得ないことがわかり，これは
\be
M_2^4 \lesssim M_{\rm Pl}^2 |\dot H| \;.
\ee
を意味する．

高次元演算子はスカラーの揺らぎだけでなくテンソルモードにも影響することに注意しよう．もっとも，これらの補正を検証するのは，おそらくはるかに難しい．例えば，スカラーの揺らぎとの関連を後で論じるユニタリーゲージの演算子 $-\bar M_3(t)/2  \cdot \delta K^\mu {}_\nu \delta
K^\nu {}_\mu $ は，$\dot g_{ij}^2$ の形の項を含む．これは重力波の分散関係を変える．実際，この演算子が存在するときのテンソルモード $\gamma_{ij}$ の作用は容易に求められ，
\be
S_{\gamma}=\frac{M_{\rm Pl}^2}{8} \int d^4x \sqrt{-g} \left[\left(1-\frac{\bar M_3^2}{M_{\rm Pl}^2}\right)\dot\gamma_{ij} \dot\gamma_{ij}- \frac1{a^2}
\partial_l\gamma_{ij} \partial_l\gamma_{ij} \right] \;.
\ee
を得る．したがって，重力波のスペクトル (\ref{eq:GWspectrum}) は $\bar M_3^2/M_{\rm Pl}^2$ 程度の補正を受ける．$\bar M_3^2$ をPlanckスケールまで大きくしない限り，この補正は小さい．$(g^{00}+1)^{n}$ の形の演算子は，計量の横波・トレースレス成分に影響しないため，テンソルモードに影響しないことは容易にわかる．

さまざまな演算子に対する実験的制限の他の例については，以下の節で論じる．

\subsectionlink{小さな音速と大きな非Gaussian性}

Goldstoneモードの作用 (\ref{Spi}) から，空間方向の運動項 $(\partial_i \pi)^2$ は背景の時間発展によって完全に決まり，$M_{\rm Pl}^2 \dot H (\partial_i\pi)^2$ となることがわかる．特に，$\dot H <0$ の場合にのみ，この項は「健全な」負の符号を持つ．これは，ヌル・エネルギー条件の破れと系の不安定性の存在との間の，よく研究されてきた関係の一例である．FRW宇宙では，この条件の破れは $\dot H<0$ と同値である \cite{Hsu:2004vr,Dubovsky:2005xd}．ただし，演算子 $(\partial_i \pi)^2$ の符号が不適切であるというだけでは，系に病的な問題があるとは結論できないことに注意しよう．以下で詳しく論じるように，特定の領域では $\delta K^\mu {}_\mu {}^2$ のような高次の項が重要になり得る．文献 \cite{Creminelli:2006xe} では，病的な問題を生じることなく $\dot H >0$ を実現できる例が研究されている．

一方，時間方向の運動項 $\dot{\pi}^2$ の係数は，二次の演算子 $(g^{00}+1)^2$ からも寄与を受けるため，背景の時間発展だけでは完全には決まらない．式~(\ref{Spi}) では，
\begin{equation}
\left(-M_{\rm Pl}^2 \dot{H} + 2 M_2^4 \right) \dot\pi^2 \;.
\end{equation}
となっている．不安定性を避けるには，$-M_{\rm Pl}^2 \dot{H} + 2 M_2^4 >0$ でなければならない．
時間方向と空間方向の運動項の係数が異なるので，$\pi$ の波の「音速」は $c_s \neq 1$ となる．背景はLorentz不変性を自発的に破っており，$c_s=1$ はいかなる対称性によっても保護されないため，これは予想されることである．前節で論じたように，$c_s=1$ からのずれは，少なくとも重力子のループによって誘起される\footnote{重力との結合と，ユニタリーゲージのラグランジアンにおける演算子の時間依存性を無視すると（そのため $\pi \to \pi + {\rm const}$ が対称性となる），$c_s=1$ は対称性 $\partial_\mu\pi \to
\partial_\mu\pi + v_\mu$ によって保護できる．ここで $v_\mu$ は定数ベクトルである．この対称変換の下で，$\pi$ のLorentz不変な運動項の変化は全微分となる一方，式~(\ref{Spi}) における $M_2^4$ に比例する演算子は明らかに不変ではないため，$c_s=1$ となる．$\pi$ により多くの微分が作用する相互作用を記述することは許されるので，この理論が自由理論ではないことに注意しよう．この対称性は，DGP模型のブレーンの曲がりモードの研究に現れる \cite{Adams:2006sv}．}．
音速は
\be
c_s^{-2} = 1-\frac{2 M_2^4}{M_{\rm Pl}^2 \dot H} \;.
\ee
で与えられる．これは，超光速伝播を避けるためには $M_2^4 >0$ でなければならないことを意味する（$\dot H <0$ を仮定する）．超光速伝播は，この理論にLorentz不変な紫外完備化が存在しないことを意味する \cite{Adams:2006sv}．以下では $c_s \leq 1$ の場合に着目する．$c_s > 1$ の模型についての現象論的な議論は，文献 \cite{Mukhanov:2005bu} を参照されたい．

上の $c_s^2$ の式を用いると，Goldstoneモードの作用は三次までで

\begin{eqnarray}\label{Spi_cs}
&&S_{\rm \pi} = \int \!  d^4 x  \:  \sqrt{- g} \left[
-\frac{M^2_{\rm Pl} \dot{H}}{c^2_s} \left(\dot\pi^2-c^2_s
\frac{(\partial_i\pi)^2}{a^2}\right)
\right.\nonumber \\ &&\qquad\left.+M_{\rm Pl}^2 \dot H \left(1-\frac{1}{c^2_s}\right)
\left(\dot{\pi}^3-\dot\pi\frac{(\partial_i\pi)^2}{a^2}
\right)- \frac43 M_3^4 \dot\pi^3... \right].
\end{eqnarray}

と書ける．

\ref{sec:Goldstone}節での議論から，エネルギー $E \gg E_{\rm mix}
\simeq M_2^2/M_{\rm Pl}$ では重力との混合を無視できることがわかっている．これは，$H$ 程度のエネルギーで行われる宇宙論的観測量の予言が，$H \gg M_2^2/M_{\rm Pl}$，あるいは同値な条件 $\epsilon/c_s^2 \ll 1$ が成り立つ場合には，Goldstoneモードの作用 (\ref{Spi_cs}) によって最低次で捉えられることを意味する．そうでない場合には，Goldstoneモードの作用が主要な効果を含んでいるという保証はない．

2点関数の計算は，$c_s=1$ の場合とほぼ同様に進められる．その際，再スケールした運動量 $\bar k=c_s k$ を用い，時間方向の運動項の前に付く追加の因子 $c_s^{-2}$ を考慮すればよい．その結果，
\be
\label{eq:spectrumzetacs}
\langle \zeta(\vec k_1) \zeta(\vec k_2)\rangle = (2 \pi)^3 \delta(\vec
k_1 + \vec k_2)  \frac1{c_{s*}} \cdot \frac{H^4_*}{4 M_{\rm Pl}^2 |\dot H_*|}\frac{1}{k_1^3} = (2
\pi)^3 \delta(\vec k_1 + \vec k_2) \frac1{c_{s*}} \cdot \frac{H_*^2}{4 \epsilon_* M_{\rm Pl}^2} \frac{1}{k_1^3} \;.
\ee
を得る．音速の時間変化は，スペクトルの傾きに追加の寄与をもたらし，
\begin{equation}
n_s=\frac{d}{d \log k}\log \frac{H^4_{*}}{|\dot{H}_*|
c_{s*}}=\frac{1}{H_*}\frac{d}{dt_*}\log \frac{H^4_{*}}{|\dot{H}_*|
c_{s*}} =4\frac{\dot{H_*}}{H^2_*}-\frac{\ddot{H}_*}{\dot{H}_*
H_* }-\frac{\dot{c}_{s*}}{c_{s*}H_{*}} \;.
\end{equation}
となる．この結果は，文献 \cite{Garriga:1999vw} で得られた結果と一致する．

作用 (\ref{Spi_cs}) から，音速を低下させるのと同じ演算子が，$\dot{\pi}(\nabla\pi)^2$ および $\dot\pi^3$ の形のGoldstoneモードの三次結合を誘起することが明確にわかる．時間微分同相変換の非線形実現は，音速の低下と3点相関関数の増大，{\em すなわち} 非Gaussian性の増大との間に関係を課す．この関係は，文献 \cite{Chen:2006nt} における3点関数の具体的な計算で強調された．

非Gaussian性の大きさを見積もるには，凍結が起こる $\omega \sim H$ の前後で，非線形補正と二次の項を比較しなければならない．
$c_s\ll 1$ の極限では，演算子 $\dot\pi(\nabla{\pi})^2$ が主要な寄与を与える．二次の作用から，モードは $k
\sim H/c_s$ で凍結することがわかり，時間微分に比べて空間微分が増強されるためである．したがって，非Gaussian性の大きさは比
\begin{equation}
\frac{{\cal L}_{\dot\pi(\nabla \pi)^2}}{{\cal L}_{2}}\sim \frac{H
\pi \left(\frac{H}{c_s} \pi\right)^2}{H^2 \pi^2} \sim \frac{H}{c^2_s}
\pi \sim \frac{1}{c^2_s} \zeta \, ,
\end{equation}
で与えられる．最後の段階では，$\pi$ と $\zeta$ の間の線形関係，式~(\ref{eq:pizeta}) を用いた．$\zeta \sim
10^{-5}$ とすると，非線形補正の大きさを見積もることができる\footnote{非線形補正の大きさは，$\zeta$ の具体的な値に依存する．$\zeta$ の典型的な値が小さくても，例えば原始ブラックホールの生成を研究するために，非常に大きな（したがって非常に起こりにくい）揺らぎに関心を持つことはあり得る．$\zeta$ の値が十分に大きく，$\zeta \gtrsim c_S^2$ となると，非線形補正は1程度となり，摂動展開は破綻する．したがって，$\zeta$ の非常に大きな値に依存する予言は，有効場理論の適用範囲外となる可能性がある．}．
通常，非Gaussian性の大きさはパラメータ $f_{\rm
NL}$ で与えられ，パラメータ依存性は ${\cal L}_{\dot\pi(\nabla
\pi)^2}/{\cal L}_2\sim f_{\rm NL} \zeta$ という形になる．したがって，主要な寄与から
\begin{equation}
\label{eq:nabla2dot}
f^{\rm equil.}_{\rm NL, \;\dot\pi(\nabla\pi)^2}\sim \frac{1}{c^2_s} \;.
\end{equation}
を得る．上付き添字の ``equil.'' は3点関数の運動量依存性を表し，これらの模型では，いわゆる正三角形型になる \cite{Babich:2004gb}．Goldstoneモードを用いた記述では，これは物理的に明らかである．関与する $\pi$ の相互作用は微分を含むため地平線の外側で急速に消失し，同程度の波長を持つモードの間にしか相関がないからである．

Goldstoneモードのラグランジアン (\ref{Spi_cs}) には，3点関数に寄与する，もう一つの独立な演算子 $-\frac43 M_3^4 \dot\pi^3$ が存在する．これは，ユニタリーゲージの演算子 $(g^{00}+1)^3$ に由来する．したがって，$\dot\pi^3$ の形の寄与は二つあり，
\begin{equation}\label{fnldotpicube}
f^{\rm equil.}_{\rm NL,\ \dot\pi^3}\sim 1-\frac{4}{3}
\frac{M_3^4}{M_{\rm Pl}^2 |\dot H| c_s^{-2}} \;.
\end{equation}
を与える．演算子 $-\frac43 M_3^4 \dot\pi^3$ の大きさは，時間微分同相変換の非線形実現によって制限されず，自由なパラメータである．DBIインフレーション \cite{Alishahiha:2004eh} では，$M_3^4 \sim M_{\rm Pl}^2 |\dot H| c_s^{-4}$ となるので，非Gaussian性への寄与は式~(\ref{eq:nabla2dot}) の寄与と同程度である．$M_3^4$ は，ユニタリーゲージの演算子 $M_2^4 (g^{00}+1)^2$ と $M_3^4 (g^{00}+1)^3$ がともに同じエネルギースケールで強結合になると仮定した場合にも，およそ同じ大きさになる．

非Gaussian性への実験的制限を，ユニタリーゲージの演算子 $(g^{00}+1)^2$ の大きさ，したがって音速に対する制約として捉えることは興味深い．具体的な計算 \cite{Chen:2006nt} によって，演算子 $\dot\pi
(\nabla\pi)^2$ が，実験的に制約されているパラメータ $f_{\rm
  NL}^{\rm equil.}$ に与える寄与が得られる．$c_s^{-1}$ について支配的な項をとると\footnote{これは，文献 \cite{Chen:2006nt} の記法で $P_{,XXX}=0$ と置くことで得られる．}
\be
f_{\rm NL}^{\rm equil.} = \frac{85}{324} \cdot \frac1{c_s^2} \;.
\ee
となる．実験的に許される範囲 \cite{Creminelli:2006rz}
\be
\label{eq:fnleqexp}
-256 < f_{\rm NL}^{\rm equil.} < 332 \quad {\rm at} \; 95\% \;{\rm C.L.}
\ee
は，制約
\be
c_s > 0.028 \quad {\rm at} \; 95\% \;{\rm C.L.}
\ee
に置き換えられる．ただし，演算子 $\dot\pi (\nabla\pi)^2$ と $\dot\pi^3$ は，原理的には3点関数に異なる運動量依存性を与えるものの，現時点ではこの差を実験的に識別できないため，制約 (\ref{eq:fnleqexp}) は二つの演算子を合わせた効果に対するものであることに注意しよう．二つの演算子の間で打ち消し合いが生じない限り，音速への制約は成り立つ．例えばDBIインフレーションでは，先ほど論じたように，演算子 $M_3^4 (g^{00}+1)^3$ の効果は無視できない．しかし，打ち消し合いは生じず，音速への制約は
\be
{\rm DBI:} \quad c_s > 0.031 \quad {\rm at} \; 95\% \;{\rm C.L.}
\ee
へとわずかに変化するだけである．

ここまではGoldstoneモードのラグランジアンに着目してきたが，この一般的な手法は，重力との混合をすべて考慮したい場合にも有用であることを強調しておきたい．例えば，ユニタリーゲージのラグランジアン (\ref{eq:actiontad}) に戻れば，3点関数の計算にいくつの係数が関与するかを容易に見て取れる．スローロールと微分の最低次では，すでに論じたとおり，係数は $M_2$ と $M_3$ の二つである．スローロールの一次では，係数の緩やかな変化を記述する新たなパラメータが四つ現れる．すなわち，通常のスローロールパラメータ $\epsilon$ と $\eta$，そして演算子 $(g^{00}+1)^2$ と $(g^{00}+1)^3$ の係数に対する二つの追加のパラメータである．実際，これは具体的な計算 \cite{Chen:2006nt} で得られる結果にほかならない\footnote{具体的な計算から，係数の一つは3点関数に独立な運動量依存性を生じないため，他のパラメータから切り分けられないことがわかる．}．

これらの議論はすべて，4点関数（およびより高次の相関関数）へ容易に拡張できる．Goldstoneモードのラグランジアンでは，4点関数に寄与する演算子が三つある（ここでもスローロールと微分の最低次で考える）．それらは，$(g^{00}+1)^2$ （音速 $c_s$ によって決まる），$(g^{00}+1)^3$，$(g^{00}+1)^4$ である．音速によって決まる演算子の効果を見積もろう．3点関数の場合と同様に，その効果は演算子 $(\nabla\pi)^4$ によって支配され，それが誘起する非Gaussian性の大きさは
\be
\frac{{\cal L}_{(\nabla \pi)^4}}{{\cal L}_{2}}\sim \frac{\left(\frac{H}{c_s} \pi\right)^4}{H^2 \pi^2} \sim \frac{H^2}{c^4_s}
\pi^2 \sim \frac{1}{c^4_s} \zeta^2 \;.
\ee
と見積もれることが容易にわかる．これは，文献 \cite{Chen:2006dfn} で行われた具体的な計算と一致する．

\subsubsectionlink{カットオフと自然性}
これまで論じたように，$c_s <1$ の場合，Goldstoneモードの作用は繰り込み不可能な相互作用を含む．したがって，Goldstoneモード間の自己相互作用は，あるエネルギースケールで強結合となり，これが本理論のカットオフを定める．このカットオフは，ツリーレベルの部分波ユニタリティを調べること，{\em すなわち} $\pi$ のツリーレベルでの散乱がユニタリである最大のエネルギーを求めることで見積もれる．計算は容易であり，唯一の複雑さは非相対論的な分散関係に由来する．カットオフスケール $\Lambda$ は，
\be
\label{eq:cutoff}
\Lambda^4 \simeq 16 \pi^2 M_2^4 \frac{c_s^7}{(1-c_s^2)^2} \simeq 16
\pi^2 M_{\rm Pl}^2 |\dot H|
\frac{c_s^5}{1-c_s^2} \;.
\ee
となる．同じ結果は，$\pi \pi$ の散乱振幅へのループ補正が重要となるエネルギースケールを調べることでも得られる．予想どおり，$c_s$ が小さくなるほど理論はより強く結合し，カットオフスケールは低くなる．一方，$c_s \to 1$ ではカットオフは次第に高くなる．この極限では繰り込み不可能な相互作用がなく，カットオフをPlanckスケールまで引き上げられるので，これは理にかなっている．このカットオフスケールは，ユニタリーゲージの演算子 $(g^{00}+1)^2$ だけを調べて得られたものである．他の独立な演算子は，その大きさによっては，さらに低いエネルギーのカットオフを与える可能性がある．スケール $\Lambda$ は，本理論が弱結合であり，意味を持つ最大のエネルギーを示していることに注意しよう．このスケールよりも下で，新しい物理が登場しなければならない．ただし，新しい物理は $\Lambda$ よりもはるかに低いスケールで現れる可能性もある．

このラグランジアンを用いて宇宙論的な相関関数を予言したいなら，インフレーション中のHubbleパラメータ $H$ 程度のスケールでこれを用いる必要がある．したがって，このエネルギースケールはカットオフよりも低く，$H \ll
\Lambda$ でなければならない．$c_s
\ll 1$ の場合にカットオフの具体的な式 (\ref{eq:cutoff}) を用いると，
\be
H^4 \ll M_{\rm Pl}^2 |\dot H| c_s^5
\ee
を得る．これは，スペクトルの規格化 (\ref{eq:spectrumzetacs}) を用いて，音速に対する不等式
\be
c_s \gg P_\zeta^{1/4} \simeq 0.003 \;. 
\ee
として書き直せる．これよりも音速が小さい理論は，$E \simeq H$ で強結合となる．当然ながら，この音速の値は，非Gaussian性が1程度となる値にも対応する．すなわち，宇宙論的予言に関係するエネルギースケール $H$ において，理論は強結合となるのである．

理論の自然性についてコメントしよう．$c_s$ が小さい極限は自然なのか，それともループ補正によって，より大きな値が誘起されるのかという疑問が生じるかもしれない．例えば $\dot\pi(\nabla\pi)^2$ や $(\nabla\pi)^4$ のようなGoldstoneモードの自己相互作用は，$(\nabla\pi)^2$ に輻射補正を誘起する．これらの寄与が $c_s^{-5}
\Lambda^4/(16 \pi^2 M_2^4)$ 程度であることは容易に見積もれる．ここで $\Lambda$ は紫外カットオフ，{\em すなわち} 新しい物理が加わるエネルギースケールである．大きな輻射補正が生じ得ないことがわかる．$\Lambda$ をユニタリティ限界 (\ref{eq:cutoff}) に取ったとしても，その効果はツリーレベルの値と同程度である．ユニタリティによるカットオフは，まさにループ補正が1程度になるエネルギースケールなので，これは理にかなっている．

また，作用 (\ref{Spi}) は有効場理論の観点から自然であることも指摘しておきたい \cite{Polchinski:1992ed}．実際，$\pi$ の近似的なシフト対称性を仮定すれば，関連演算子は大きな繰り込みから保護される．この場合，関連演算子の係数は十分に小さく，観測にとって重要になることはない．宇宙論的な相関関数が探るのは，$H$ 程度の固定されたエネルギースケールでの理論であり，それよりも低いエネルギーへ進むことはないからである．もちろん，ここではインフレーションの期間だけを考えており，この期間には，近似的なシフト対称性があれば理論を技術的に自然にするのに十分である．インフレーションからの円滑な終了と効率的な再加熱を実現することは，機能する模型に対する追加の要件であり，本形式では論じない．

\subsectionlink{de Sitter極限とゴースト凝縮}

前節では，$c_s \to 0$ の極限では理論がますます強結合になるため，この極限が病的であることを見た．しかし，これまでの議論では，ユニタリーゲージのラグランジアン（\ref{eq:actiontad}）に含まれる高階微分演算子
\be
\int d^4 x \sqrt{-g} \left(-\frac{\bar M_2(t)^2}{2} \delta K^\mu {}_\mu {}^2
-\frac{\bar M_3(t)^2}{2} \delta K^\mu {}_\nu \delta K^\nu {}_\mu
\right) \;.
\ee
を無視していた．これらの演算子は，Goldstone作用において
\begin{equation}
\label{di_pi2}  \int  \! d^4x \: \sqrt{-g} \left[ - \frac{\bar M^2}{2} \, \frac{1}{a^4}(\di_i^2 \pi)^2
   \right] \; ,
\end{equation}
という形の空間的運動項を生む．ここで，$\bar M^2 = \bar M_2^2+  \bar M_3^2$ である．
この空間的運動項によって，$c_s \to 0$ の極限でもGoldstoneモードは伝播する．したがって，$\dot H = 0$ の極限で，一般的なラグランジアンを考察することは興味深い．この極限では重力背景は厳密なde Sitter時空となり，$c_s=0$ が成り立つ．この場合，$H$ は時間に依存しないため，理論に追加の対称性を課すことができる．すなわち，ユニタリーゲージのラグランジアンに含まれるすべての係数を時間に依存しないものとする．Goldstoneモード $\pi$ を再導入する手順（\ref{Smixed}）を振り返ると，この対称性は微分を伴わない $\pi$ への依存をすべて禁止することがわかる．したがって，Goldstone作用は $\pi$ のシフト
\be
\pi(\vec x,t) \to \pi(\vec x,t) + {\rm const.}
\ee
のもとで不変である．これがゴースト凝縮 \cite{Arkani-Hamed:2003pdi} の極限であり，Goldstoneモードは非相対論的な分散関係 $\omega \propto k^2$ を持つ．より一般には，空間的運動項 $c_s^2
(\nabla\pi)^2$ と高階微分項 $(\nabla^2\pi)^2$ の両方が存在する中間的な状況を考えることができる．理論の予言は，凍結が起こるエネルギー $\omega
\sim H$ においてどちらの項が支配的であるかによって，大きく変化する \cite{Arkani-Hamed:2003juy,Senatore:2004rj}．

これまでの模型と同様に，Goldstoneモードと重力の混合を無視できるエネルギー領域を求めなければならない．この領域を見積もる簡単な方法の一つは，$\delta K$ 演算子が含む次のような項を見ることである．
\be
\delta K_{ij} \supset (\partial_i\partial_j \pi +\partial_i g_{0j}) \;.
\ee
正準規格化を行うと，$k \gtrsim M_2^2/M_{\rm Pl}$ では重力との混合を無視できることがわかる．Goldstoneモードの分散関係は $\omega^2 = \bar M^2/M_2^4 \cdot k^4$ の形なので，混合が重要になるエネルギーの上限 $E_{\rm mix}$ は
$E_{\rm mix} \simeq \bar M M_2^2/M_{\rm Pl}^2$ であることがわかる \cite{Arkani-Hamed:2003pdi}．
このスケールは背景の曲率とは無関係であることに注意しよう．これは，この例の非常に顕著な特徴である．通常，重力との混合は背景のエネルギー・運動量テンソル，したがって時空の曲率に関係しており，系が空間を強く湾曲させるほど，重力との混合も強くなる．一方，この場合には，平坦なMinkowski背景上でも混合が重要になる．これが本来の意味での重力の修正と呼ばれるものである．すなわち，たとえば源が生成するNewtonポテンシャルのような重力が，曲率スケールよりもはるかに小さなスケールで修正される．この修正重力模型とその豊かな現象論は，文献 \cite{Arkani-Hamed:2003pdi} で研究されている\footnote{音速が低下した模型の場合にも，重力との混合のスケールは地平線よりパラメトリックにはるかに小さくなりうる．$\epsilon/c_s^2 \ll 1$ であれば十分である．この場合，その模型は重力を修正する一つの方法と見なすことができる．ただし，空間的な高階微分項を考慮せずに $\dot H = 0$ の極限をとることはできない．そうしなければ，スカラー・モードは伝播しなくなる．}．

ここではインフレーションに関心があるので，逆の極限 $H \gg E_{\rm mix}$ に注目する．この極限では混合を無視でき，$\pi$ のラグランジアンに集中できる．ゴースト・インフレーションの主な特徴を簡単に説明しよう．詳細については文献
\cite{Arkani-Hamed:2003pdi,Arkani-Hamed:2003juy,Senatore:2004rj} を参照されたい．これらの文献では，本論文で提示したものに非常に近いアプローチで理論が研究されている．興味深い特徴の大半は，各演算子のエネルギーに関するスケーリングを見ることで理解できる．非相対論的な分散関係 $\omega \propto k^2$ のため，演算子のエネルギーに関するスケーリングは，Lorentz不変な場合のようにその質量次元とは一致しない．
エネルギーを因子 $s$ で再スケールする変換 $E\rightarrow s E$（時間の再スケール $t\rightarrow s^{-1} t$ と等価である）には，運動量の変換 $k\rightarrow s^{1/2}k$（空間座標では $x\rightarrow
s^{-1/2}x$）が伴わなければならない．Goldstoneモードの二次の作用は
\be
\label{eq:kinghost}
\int d^4 x \left[ 2 M_2^4 \dot\pi^2 - \frac{\bar M^2}{2} \,
  \frac{1}{a^4}(\di_i^2 \pi)^2 \right]
\ee
という形なので，二次の作用を不変に保つには，$\pi$ にスケーリング次元 $1/4$ を割り当てなければならない．すなわち，
\begin{equation}
\pi\rightarrow s^{1/4}\pi \, .
\end{equation}
この規則を用いると，運動項（\ref{eq:kinghost}）以外の，許されるすべてのGoldstone演算子が非関連であることを容易に確かめられる．{\em すなわち}，これらは正のスケーリング次元を持ち，エネルギーを下げるほど重要性を失っていく．このことは，この理論が有効場理論として意味を持つことを示している．特に，高階時間微分を含む運動演算子 $\ddot \pi^2$ は，一見すると $(\nabla^2\pi)^2$ と同程度に重要に見え，理論にゴーストが存在することを記述するが，次元が 2 であり，低エネルギーでは無視できる．問題に単一のスケール $M$ があり，$M \simeq M_2 \simeq \bar M$ であると仮定すれば，このスケールが有効場理論による記述のエネルギー・カットオフも定める．

$\pi$ のスケーリング次元からは，ゴースト・インフレーションで生成される摂動のスペクトルも見積もることができる．$\pi$ の次元は，量子揺らぎの振幅がエネルギーとともにどのように変化するかを教えてくれる．カットオフのスケール $\Lambda \simeq M$ では，正準規格化されたGoldstone場 $\pi_c \simeq M^2 \pi $ の量子揺らぎは，カットオフと同程度の大きさ $\delta \pi_c (M)\sim M$ である．エネルギーを下げることで，凍結時 $E \sim H$ の量子揺らぎを見積もることができる．標準的な場合，スカラー場の次元は 1 であり，凍結時の揺らぎは $H$ 程度となるが，ここでは
\begin{equation}
\delta\pi_c(H) \sim \delta\pi_c(M)
\left(\frac{H}{M}\right)^{1/4}\sim \left(H M^3\right)^{1/4} \;.
\end{equation}
となる．スケール $H$ での量子揺らぎは，Lorentz不変な分散関係を持つスカラー場と比べて大幅に増大している．したがって，$\zeta$ のスペクトルは \cite{Arkani-Hamed:2003juy}
\begingroup\addtolength{\belowdisplayskip}{-2pt}
\be
\label{eq:spectrumzetaghost}
\langle \zeta(\vec k_1) \zeta(\vec k_2)\rangle \sim (2 \pi)^3 \delta(\vec
k_1 + \vec k_2)  \frac{H^2}{M^4} (H M^3)^{1/2} \frac{1}{k_1^3} \;.
\ee
\endgroup
で与えられる．スペクトルが正しい規格化を持つためには，$(H/M)^{5/4} \simeq
10^{-5}$ が必要である．

非Gaussian性は，次元が最も低い演算子によって支配される．演算子 $M_2^4
\dot\pi (\nabla\pi)^2$ は，ユニタリーゲージの演算子 $M_2^4
(g^{00}+1)^2$ から生じ，次元 $1/4$ を持つ最も非関連性の弱い演算子であることが容易にわかる．
カットオフのスケール $M$ では，理論は強結合である．三次の演算子は次元 $1/4$ を持つため，$H$ 程度のエネルギーでは，$(H/M)^{1/4}$ 程度の非Gaussian性を生み，これはパラメトリックには $P_\zeta^{1/5}$ 程度である．同じ結果は，前節で用いたアプローチ，{\em すなわち}凍結時に相互作用項と自由作用を比較することによっても得られる．
\begin{equation}
\label{eq:NGghost}
\frac{{\cal L}_{\dot\pi(\nabla \pi)^2}}{{\cal L}_{2}}\sim
\frac{M^4 H \pi (H M) \pi^2}{M^4 H^2 \pi^2}=M \pi=\frac{M}{H}
\zeta\sim\left(\frac{H}{M}\right)^{1/4} \, .
\end{equation}
最も関連性の高い演算子の次元が非常に低いため，非Gaussian性の大きさは標準的なスローロールの場合と比べて極めて大きい．具体的な計算
\cite{Arkani-Hamed:2003juy} では，効果は素朴な見積もりよりいくぶん小さくなり，既存の実験的な制限 \cite{Creminelli:2006rz} と同程度になる．

これまでの議論では，ユニタリーゲージの演算子
\begin{equation}
\label{eq:leooperator}
\int \! d^4 x \: \sqrt{- g} \:
\left(-\frac{\bar M_1(t)^3}{2} (g^{00}+1)\delta K^\mu {}_\mu \right)\;.
\end{equation}
を無視してきた．この演算子は時間反転のもとで奇であるため，その係数をゼロにすることは無矛盾である．この項が存在すれば，Goldstoneラグランジアンには，次元 $1/4$ を持つ二つ目の演算子が，$\nabla^2\pi (\nabla\pi)^2$ という形で現れる．その 3 点関数への寄与は，$\dot\pi(\nabla\pi)^2$ の寄与と同程度となる．ユニタリーゲージの演算子（\ref{eq:leooperator}）は，これから論じるように，$\pi$ の二次のラグランジアンにも寄与する．

\subsubsectionlink{ゴースト凝縮を伴わないde Sitter極限}
本節では，演算子（\ref{eq:leooperator}）が $\pi$ の二次の作用に及ぼす効果を調べる．この演算子の係数が十分に大きければ，新しいde Sitter極限が得られることを見る．この極限では，凍結時の分散関係が $\omega^2
\propto k^2$ の形となり，ゴースト凝縮の振る舞い $\omega^2
\propto k^4$ に取って代わる．

簡単のため，$\bar M_1$ は時間に依存しないとする．Goldstoneモードを再導入すると，$-\bar M_1^3 \dot\pi\nabla^2\pi/a^2$ という形の三階微分項が得られる（\footnote{この演算子は，$\dot\pi^2$ に対して，$H$ に比例する寄与も与える．この寄与は $M_2^4 \dot\pi^2$ に比べて小さいと仮定する．Minkowski時空では，ここで考えている演算子を $\phi \to -\phi$ 対称性によって禁止できる．この対称性は，ユニタリーゲージでは時間反転と等価である \cite{Arkani-Hamed:2003pdi}．de Sitter背景では，この対称性は計量によって破られるため，この演算子の係数をゼロにすることはできない．}）．スケール因子 $a(t)$ と計量の行列式に時間依存性がなければ，これは時間に関する全微分となる．部分積分すると，標準的な二階微分の空間的運動項
\be
\label{eq:byparts}
-\int d^4 x \sqrt{-g} \,\frac{\bar M_1^3 H}{2}
\left(\frac{\partial_i}{a}\pi\right)^2 \;.
\ee
が得られる．厳密なde Sitter極限 $\dot H =0$ で，$M_2 \sim \bar
M_1 \sim M$ とすると，この演算子は $\omega^2 = c_s^2 k^2$ という形の分散関係を与え，音速は小さくなる\footnote{この模型では，重力との混合はこれまでの場合とかなり異なる．その理由は，式~（\ref{eq:byparts}）を得るために時間微分を部分積分しているため，Goldstoneの項が，重力との混合を記述する項と同じ数の微分を含むことにある．このことは，高エネルギーに向かっても混合の重要性が低下していかないことを意味する．一方，混合が常に無視できるように模型のパラメータを選ぶことは可能である．具体的な計算については文献 \cite{Creminelli:2006rz} を参照されたい．}．
\begin{equation}
c^2_s=\frac{H}{M}\ll 1 \;.
\end{equation}
これが成り立つのは，高階微分演算子 $\delta K^\mu {}_\mu {}^2$ と
$\delta K^\mu {}_\nu \delta K^\nu {}_\mu$ が劣勢である場合に限られる．これらも同じ質量スケールで特徴づけられ，$\bar M_2 \sim \bar
M_3 \sim M$ であると仮定すると，分散関係には二つの寄与が生じる．
\be
\omega^2 \sim \frac{H}{M} k^2 + \frac{k^4}{M^2} \;.
\ee
この二つの空間的運動項は，凍結時 $\omega \sim
H$ には同程度となる．一方，$k^4$ の寄与がいくらか抑制されていれば，凍結時には無視でき，したがってインフレーションの予言にとっても重要ではなくなる．この極限では，厳密なde Sitter背景を持つ一方で，凍結時には $\omega^2 \propto k^2$ の分散関係を持つ，新しい種類のゴースト・インフレーションが得られる．

他の模型に対して行った手順に従えば，スペクトルの規格化と，3 点関数の非Gaussian性の見積もりを容易に得ることができる．
\be
\label{eq:spectrumzetaleo}
\langle \zeta(\vec k_1) \zeta(\vec k_2)\rangle \sim (2 \pi)^3 \delta(\vec
k_1 + \vec k_2)  \frac{H^4}{M^4} \frac1{c_s^3} \frac{1}{k_1^3}  \sim (2 \pi)^3 \delta(\vec
k_1 + \vec k_2)  \left(\frac{H}{M}\right)^{5/2} \frac{1}{k_1^3} \;.
\ee
\be
\frac{{\cal L}_{\dot\pi(\nabla \pi)^2}}{{\cal L}_{2}}\sim
\frac1{c_s^2} H \pi \sim\left(\frac{H}{M}\right)^{1/4} \, .
\ee
Goldstone演算子 $(\nabla\pi)^2\nabla^2\pi$ からも，同程度の寄与が生じる．当然ながら，上の見積もりは，ゴースト凝縮の場合に得た式（\ref{eq:spectrumzetaghost}）および（\ref{eq:NGghost}）と同じである．実際，すでに論じたように，すべての演算子を同じスケールにとると，凍結時には $k^2$ と $k^4$ の空間的運動項から同程度の寄与が得られる．したがって，二つの寄与のうち一方だけが存在すると仮定した場合にも，同様の予言が得られると期待される．

凍結時に $(g^{00}+1)\delta K^\mu {}_\mu $ が支配する場合と，$\delta K^\mu {}_\mu {}^2$ および $\delta K^\mu {}_\nu \delta K^\nu
{}_\mu$ が支配する場合という，二つの異なるde Sitter極限が見つかった．そこで，他の可能性もあるのではないかと考えるかもしれない．何らかの理由で，これら二つの空間的運動項がともに抑制され，より高次で初めて主要な演算子が現れることを想像できる．この場合，最終的には
\begin{equation}
\omega^2\sim k^{2n} \quad n \geq 3 \;.
\end{equation}
という形の分散関係が得られる．しかし，このようなことは起こりえないと容易にわかる．なぜなら，その理論は有効場理論として意味を持たなくなるからである．ゴースト凝縮で用いたのと同じ論理に従うと，演算子 $\pi$ のスケーリング次元は
\begin{equation}
\pi\rightarrow s^{-\frac12+\frac{3}{2n}} \pi \, .
\end{equation}
となる．このことから，演算子 $\dot\pi(\nabla\pi)^2$ は，対称性によって時間的運動項 $\dot\pi^2$ と結びついており，次元
$(7-3n)/(2n)$ を持つ．$n \geq 3$ では，この演算子は低エネルギーで強結合となるため，有効場理論は意味を持たない．

\sectionlink{結論}
インフレーションを検証するための実験的な取り組みが進行し，さまざまな模型が数多く提案されていることを考えると，インフレーションの最も一般的な理論を特徴づけることは非常に重要である．本論文では，新しい観点を採用した．一般的なラグランジアンを書き下し，インフレーションを起こす解のまわりの揺らぎを調べる代わりに，準de Sitter背景のまわりの揺らぎの有効理論を直接記述する．この理論では，空間微分同相変換は顕在的であり，時間微分同相変換は非線形に実現される．微分の最低次では，最も一般的な作用を次の形で書けることを示した．
\be
S &=& \int  \! d^4 x \; \sqrt{- g} \Big[ \frac12 M_{\rm Pl}^2 R + M_{\rm Pl}^2 \dot H
g^{00} - M_{\rm Pl}^2 (3 H^2 + \dot H) +
 \nonumber \\ && \frac{1}{2!}M_2(t)^4(g^{00}+1)^2+\frac{1}{3!}M_3(t)^4 (g^{00}+1)^3+
 \nonumber \\
&& - \frac{\bar M_1(t)^3}{2} (g^{00}+1)\delta K^\mu {}_\mu
-\frac{\bar M_2(t)^2}{2} \delta K^\mu {}_\mu {}^2
-\frac{\bar M_3(t)^2}{2} \delta K^\mu {}_\nu \delta K^\nu {}_\mu + ...
\Big] \; .
\ee

宇宙論的な相関関数は，Hubbleパラメータ $H$ 程度のスケールでこの有効場理論を検証する．このアプローチでは，対称性の役割がはるかに明確になる．どの特徴がインフレーションを起こす背景解から導かれるかを具体的に見ることができる．特に，$\dot H<0$，$\dot H = 0$，$\dot H > 0$ の各場合で，振る舞いが大きく異なることがわかる．これは，演算子 $g^{00}$ の係数が $\dot H$ によって固定されるためである．この観点から，本アプローチは，インフレーション，修正重力理論，およびバウンス・模型 \cite{Creminelli:2006xe,Buchbinder:2007ad,Creminelli:2007aq} のようにヌル・エネルギー条件を破る理論（宇宙論の文脈では $\dot H >0$ と等価）相互の関係をより明確にする．対称性の役割を示すもう一つの例は，音速の低下と非Gaussian性の増大の結びつきである．両者は同じ演算子 $M_2(t)^4(g^{00}+1)^2$ から生じるため，時間微分同相変換の非線形実現によって関連づけられている．

最も標準的なスローロールのシナリオからのあらゆる逸脱は，素粒子物理の標準模型の研究におけるのと同様に，高次演算子の大きさに体系的に組み込まれている．さらに，すべての単一場模型が共通の枠組みに統一されるため，たとえば文献 \cite{Cheung:2007sv,Arkani-Hamed:2007ryv} で行われたように，具体的な実現の仕方に依存しない一般的な結論を導くことが可能になる．

本定式化の拡張として考えられるものは容易に思い浮かぶ．同じ方針に沿って，（単一場の）クインテッセンスの最も一般的な理論を調べ，その現象論的な帰結を求めることは興味深いだろう．有効理論を $H$ 程度のスケールで調べるインフレーションとは異なり，この場合に関心があるのは，地平線内での摂動の力学である．放射優勢期や物質優勢期のような流体における揺らぎの研究 \cite{Dubovsky:2005xd} に，本アプローチを用いることも興味深い．最後に，追加の場を導入して，多場インフレーション・模型を研究することも容易であるはずだ．

\phantomsection\section*{\hyperlink{toc}{謝辞}}\addcontentsline{toc}{section}{謝辞}
多くの有益な議論をしていただいた N.~Arkani-Hamed，L.~Boubekeur，S.~Dubovsky，
A.~Guth，F.~Vernizzi，J.~Wacker，M.~Zaldarriaga，そして特に A.~Nicolis に感謝する．
A. Liam Fitzpatrick は NSF，Jared Kaplan は Hertz と NSF のフェローシップの支援を受けている．

\appendix

\section*{\hyperlink{toc}{付録}}

\sectionlink{ユニタリーゲージにおける最も一般的なラグランジアン\label{app:generic}}
ユニタリーゲージにおける最も一般的なラグランジアンを書き下すための規則を調べよう．空間微分同相変換のもとでのみ不変な理論には，関数 $\tilde t(x)$（時間的な勾配を持つ）が定める特別な時空の葉層分解があり，この関数が時間微分同相変換を非線形に実現する．
たとえば，時間発展するスカラー場によって対称性が破られている場合，$\tilde t$ が一定の面は，スカラー場の値が一定の面でもある．ユニタリーゲージとは，時間座標 $t$ が $\tilde t$ と一致するように選んだゲージであり，このため追加の自由度 $\tilde t$ は作用に陽には現れない．したがって，次のようなさまざまな項を構成できる．
\begin{enumerate}
\item すべての微分同相変換のもとで不変な項．これらは，Riemannテンソル $R_{\mu\nu\rho\sigma}$ とその共変微分からなる多項式を，スカラーになるように縮約したものにほかならない\footnote{添字の縮約には，計量と完全反対称テンソル $(-g)^{-1/2} \epsilon^{\mu\nu\rho\sigma}$ を用いることができる．}．
\item $\tilde t$ の一般的な関数は，ユニタリーゲージでは $f(t)$ になる．したがって，ラグランジアンの任意の項の前に，時間の一般的な関数を自由に掛けることができる．
\item 勾配 $\partial_\mu \tilde t$ は，ユニタリーゲージでは $\delta_\mu^0$ になる．したがって，すべてのテンソルにおいて，上付きの $0$ 添字を縮約せずに残すことが常にできる．たとえば，ユニタリーゲージのラグランジアンでは，$g^{00}$（およびその関数）や，Ricciテンソルの成分 $R^{00}$ を用いることができる．
\item $\tilde t$ が一定の面に垂直な単位ベクトルを定義すると便利である．
\be
n_\mu = \frac{\partial_\mu\tilde t}{\sqrt{-g^{\mu\nu}\partial_\mu \tilde t \partial_\nu\tilde t}}  \;.
\ee
これにより，$\tilde t$ が一定の面上の誘導空間計量 $h_{\mu\nu} \equiv g_{\mu\nu} + n_\mu n_\nu$ を定義できる．あらゆるテンソルは，$h_{\mu\nu}$ を用いてこれらの面上に射影できる．特に，誘導された三次元計量のRiemannテンソル $^{(3)} R_{\alpha\beta\gamma\delta}$ と，三次元計量に関する共変微分を作用に用いることができる．
\item $\partial_\mu\tilde t$ の共変微分から，さらに別の可能性が得られる．同じこととして，$n_\mu$ の共変微分を考えてもよいことに注意しよう．規格化因子に微分が作用しても，$\partial_\mu g^{00}$ のような項が得られるだけであり，これらはそれ自体で共変で，ユニタリーゲージのラグランジアンに用いることができる．$n_\mu$ の共変微分を $\tilde t$ が一定の面上に射影すると，これらの面の外的曲率が得られる．
\be
K_{\mu\nu} \equiv h_\mu^\sigma \nabla_\sigma n_\nu \;.
\ee
添字 $\nu$ はすでに面上に射影されている．これは，$n^\nu \nabla_\sigma n_\nu = \frac12 \nabla_\sigma (n^\nu n_\nu) =0$ であるためである．面に垂直な方向の $n_\nu $ の共変微分は，
\be
n^\sigma \nabla_\sigma n_\nu = - \frac12 (-g^{00})^{-1} h^\mu_\nu \partial_\mu(-g^{00})
\ee
と書き換えられるため，新しい項を生じない．したがって，$n_\mu$ のすべての共変微分は，外的曲率 $K_{\mu\nu}$（およびその共変微分）と $g^{00}$ の微分を用いて書くことができる．

\item 誘導された三次元計量のRiemannテンソルと外的曲率を同時に用いるのは冗長であることに注意しよう．$^{(3)} R_{\alpha\beta\gamma\delta}$ は，Gauss--Codazzi関係を用いて \cite{Wald:1984rg}
\be
^{(3)} R_{\alpha\beta\gamma\delta}= h^\mu_\alpha h^\nu_\beta h^\rho_\gamma h^\sigma_\delta R_{\mu\nu\rho\sigma} -K_{\alpha\gamma}K_{\beta\delta} +K_{\beta\gamma}K_{\alpha\delta} \;.
\ee
と書き換えられるからである．したがって，三次元Riemannテンソルを用いる必要はまったくない．誘導計量 $h_{\alpha\beta}$ を陽に用いることも避けられる．四次元計量と $n_\mu$ で書き表せば，すでに上で論じた項だけが得られる\footnote{誘導計量の行列式は，四次元の全計量の行列式と $h = g^{00} \cdot g$ で関係づけられること，また，三次元の完全反対称テンソルは，四次元のものを用いて $h^{-1/2} \epsilon^{ijk} = (-g)^{-1/2} (-g^{00})^{-1/2} \epsilon^{0ijk}$ と書き換えられることに注意しよう．}．最後に，誘導された三次元計量に関する共変微分を用いることも避けられる．射影されたテンソルの三次元共変微分は，四次元共変微分を射影することで得られる \cite{Wald:1984rg}．
\end{enumerate}

以上より，ユニタリーゲージにおける最も一般的な作用は
\be
S = \int d^4 x \sqrt{-g} \; F(R_{\mu\nu\rho\sigma}, g^{00}, K_{\mu\nu},\nabla_\mu,t) \;,
\ee
で与えられると結論できる．ここで，関数 $F$ の内部で縮約されずに残る添字は，すべて上付きの 0 でなければならない．

\sectionlink{与えられた FRW 解のまわりでの展開\label{app:tadpoles}}
本節では，与えられた FRW 背景（空間曲率に応じて $k=-1,0,+1$）のまわりで，時間微分同相変換が破れた最も一般的な理論を，
\be
S = \int d^4x \sqrt{-g} \left[\frac12 M_{\rm Pl}^2 R + M_{\rm Pl}^2 \Big(\dot H-\frac{k}{a^2}\Big) \cdot g^{00} - M_{\rm Pl}^2 \Big(3 H^2 +\dot H+2\frac{k}{a^2}\Big) + \ldots \right]
\label{eq:general}
\ee
と書けることを証明したい．ここで，省略記号は，空間微分同相変換のもとで不変であり，与えられた FRW 背景のまわりの揺らぎについて二次（またはそれ以上）の項を表す\footnote{$R$ の前の係数は，適切な場の再定義 $g_{\mu\nu} \to g_{\mu\nu} \cdot
  f(t)$ によって，常に時間に依存しないものにできる．これは通常の定式化では，Einsteinフレームへ移ることに対応する．}．

そのままでは，この主張は自明である．明示した項が望む FRW 発展を生むことはわかっているので，そこから離れないようにするには，追加の演算子はこの解のまわりで二次から始まらなければならない．ここで主張したいのは，追加される不変な項の{\em 一つ一つ}が，異なる演算子の間で一次の寄与を相殺させることなく，二次（またはそれ以上）となるということである．これらの項は，$g^{00}+1$，$\delta K_{\mu\nu}= K_{\mu\nu} -K_{\mu\nu}^{(0)}$，$\delta R_{\mu\nu\rho\sigma} \equiv R_{\mu\nu\rho\sigma} - R_{\mu\nu\rho\sigma}^{(0)}$ などの一次の演算子からなる，多項式（二次以上）として書かれる．与えられた FRW 解で評価した値を陽に差し引いているため，これらの演算子は摂動の一次から始まることに注意しよう．FRW 計量の対称性により，背景上で評価したあらゆるテンソル（$K_{\mu\nu}^{(0)}$，$R_{\mu\nu\rho\sigma}^{(0)}$，$(\nabla_\alpha R_{\mu\nu\rho\sigma})^{(0)}$ など）は，$g_{\mu\nu}$，$n_\mu$，および時間の関数だけを用いて書ける．たとえば，
\begin{eqnarray}
K^{(0)}_{\mu\nu} & =&  a^2 H h_{\mu\nu} \\
R_{\mu\nu\rho\sigma}^{(0)} & = & 2 (H+k) h_{\mu [\rho} h_{\sigma ] \nu} + (\dot H + H^2) a^2 h_{\mu\sigma} \delta_\nu^0 \delta_\rho^0 + {\rm perm.} 
\end{eqnarray}
となる．ここで，$k$ は空間的切片の曲率に依存する定数であり，置換は最後の項にのみ作用する．このように，FRW 背景上で評価したすべての演算子は，それ自体が共変な演算子である．したがって，$\delta K_{\mu\nu}$ や $\delta R_{\mu\nu\rho\sigma}$ のような演算子は，適切に定義された共変な演算子であり，与えられた FRW 背景上ではゼロとなり，摂動の一次から始まる．
背景上で評価したテンソルをこのように書き換えられるのは，FRW 背景の高い対称性のおかげに限られることを強調しておく．対称性の低い解，{\em たとえば}非一様な背景のまわりで展開する場合には，これは成り立たない．

式（\ref{eq:general}）では，明示した演算子だけが揺らぎの一次の項を含むため，$\sqrt{-g}\;g^{00}$ と $\sqrt{-g}$ の係数は，背景の FRW 解によって一意に決定される．

ラグランジアンを常に（\ref{eq:general}）の形にできることを見てみよう．二つのテンソル $T$ と $G$ の縮約からなる演算子をとると（より多くのテンソルへの一般化は容易である），
\be
TG = \delta T \delta G + T^{(0)} G + T G^{(0)} - T^{(0)} G^{(0)} \;. \label{eq:splitting}
\ee
と書ける．和の各項について考えよう．第 1 項は，望みどおり，摂動の二次から始まることが明示されている．すでに述べたように，FRW 背景の対称性により，非摂動のテンソル $T^{(0)}$ と $G^{(0)}$ は，$g_{\mu\nu}$，$n_\mu$，$t$ の関数として書ける．したがって，最後の項 $T^{(0)} G^{(0)}$ は，時間に依存する係数を持つ $g^{00}$ の多項式にすぎない．これは，$\sqrt{-g}\;g^{00}$ と $\sqrt{-g}$ の項に加えて，摂動の二次から始まることが明示された演算子を含む．
残るのは，$T^{(0)} G$ という形のテンソルである．これらの項も，式~（\ref{eq:general}）の一次の演算子と，揺らぎの二次から始まる演算子との和として書けることを証明したい．構成上，$G$ は $K_{\mu\nu}$ または $R_{\mu\nu\rho\sigma}$ のいずれかについて一次であり，それらに共変微分が作用したものである．共変微分は，繰り返し部分積分を行い，$T^{(0)}$ と演算子の時間に依存する係数に作用させることで処理できる．その際，外的曲率の項が生成されることがある．その場合は，共変微分がなくなるまで式~（\ref{eq:splitting}）を繰り返し適用すればよい\footnote{$g^{00}$ のべきが，$T^{(0)}$ から生じることもある．これらは $g^{00} = -1 +\delta g^{00}$ と書くことで処理でき，式~（\ref{eq:general}）の $g^{00}$ 演算子への追加の寄与と，摂動について陽に二次以上である項とが生成される．}．この結果，共変微分を含まない，可能なスカラーの一次の項だけが残る．すなわち，$K^{\mu}  {}_{\mu}$ と $R^{00}$ である．どちらも，より便利な形に書き換えられる．外的曲率の項は部分積分でき，
\begin{eqnarray}
&&\int d^4x \sqrt{-g}\; f(t) K^{\mu} {}_{\mu} =\int d^4x \sqrt{-g}\; f\; \nabla_{\mu} n^{\mu}\nonumber \\ &&\qquad=-\int d^4 x \sqrt{-g} \;n^\mu \partial_\mu f =\int d^4x \sqrt{-g} \sqrt{-g^{00}} \dot f \,. 
\end{eqnarray}
となる．一方，$R^{00}$ は，次の関係式 \cite{Wald:1984rg} を用いて処理できる．
\begin{eqnarray}
(-g^{00})^{-1} R^{00}=R_{\mu\nu}n^\mu n^\nu = K^2 -K_{\mu\nu}K^{\mu\nu}-\nabla_\mu(n^\mu\nabla_\nu n^\nu)+\nabla_\nu(n^\mu \nabla_\mu n^\nu)\, .
\end{eqnarray}
最後の二つの項は，再び部分積分できる．
\begin{eqnarray}
\int d^4x \sqrt{-g} f(t) \nabla_\mu(n^\mu\nabla_\nu n^\nu)=-\int d^4x \sqrt{-g}\; \partial_\mu f\; n^\mu K^\nu {}_\nu\, ,
\end{eqnarray}
\begin{eqnarray}
\int d^4 x \sqrt{-g} f(t) \nabla_\nu(n^\mu \nabla_\mu n^\nu)=-\int d^4 x \sqrt{-g}\; \partial_\nu f\; n^\mu\nabla_\mu n^\nu =0
\end{eqnarray}
最後の変形では，$\partial_\nu f \propto n_\nu$ を用いた．以上から，$K^\mu_\mu$ と $R^{00}$ は，式~（\ref{eq:general}）の一次の演算子と，揺らぎの二次から始まる不変な項とを用いて書けることが示された．

以上より，与えられた FRW 背景のまわりで時間微分同相変換が破れた理論の最も一般的なラグランジアンは，次の形に書けることを示した．
\begin{eqnarray}\label{eq:general2}
S &=& \int d^4x \sqrt{-g} \; \Big[\frac12 M_{\rm Pl}^2 R + M_{\rm Pl}^2 \Big(\dot H-\frac{k}{a^2}\Big) \cdot g^{00} - M_{\rm Pl}^2 \Big(3 H^2 +\dot H+2\frac{k}{a^2}\Big) +  \\ \nonumber &&  F^{(2)}(g^{00}+1,\delta K_{\mu\nu}, \delta R_{\mu\nu\rho\sigma};\nabla_\mu;t)\Big]
\end{eqnarray}
ここで，$F^{(2)}$ は，引数 $g^{00}+1$，$\delta K_{\mu\nu}$，$\delta R_{\mu\nu\rho\sigma}$ について二次から始まる．



\clearpage
\renewcommand{\refname}{\hyperlink{toc}{参考文献}}
\bibliographystyle{utphys}
\bibliography{ref}
\clearpage
\phantomsection
\section*{\hyperlink{toc}{訳注：原文の表記について}}
\addcontentsline{toc}{section}{訳注：原文の表記について}
本訳は指定されたarXiv第2版の全文を対象とする．本文中の数式，式番号，引用番号を原文どおり保持し，参考文献の題名と書誌情報は原語で収録した．以下は原文との照合で見つかった確認事項であり，原著者による訂正ではない．
\begin{itemize}
\item 第4.2節冒頭は，ヌル・エネルギー条件の破れに対応する条件を$\dot H<0$と記している．直後の議論および$\rho+p=-2M_{\rm Pl}^2\dot H$との関係からは，$\dot H>0$の誤記と考えられるが，本文は原文の符号を保持した．
\item 式(41)左辺は原文どおり$n_s$とした．式(33)で用いられる定義との比較からは，$n_s-1$とすべき可能性がある．
\item 第2節の$\delta K_{\mu\nu}$の定義および式(77)には$a^2Hh_{\mu\nu}$が現れる．誘導計量$h_{\mu\nu}$の定義との整合性に疑義があるが，原文を保持した．式(78)の第1項も，原文PDFとTeX原稿の双方にある$2(H+k)h_{\mu[\rho}h_{\sigma]\nu}$のままである．
\item 第3節の正準規格化の説明には$\dot H^{1/2}$という表記がある．通常のスローロールでは$\dot H<0$なので，規格化の大きさについては$|\dot H|^{1/2}$と読む必要があるが，本文の式は変更していない．
\item 第4.1節の式(36)直前の本文には，$\bar M_3(t)$を係数とする表記がある．式(10)および式(36)の$\bar M_3^2$との違いを含め，原文どおり収録した．
\item 第3節末尾の英文は，非Gaussian性についての例示の係り受けが曖昧である．本訳では第4.1節の明示的な説明に従い，通常のスローロール・インフレーションの非Gaussian性を，重力との混合が重要な例として訳した．
\item 参考文献はINSPIRE-HEPから取得したBibTeXに基づき、原文の引用番号と対応順を維持して収録した。原文では文献[16]が、普遍的補正に関する最近の議論として紹介されている。書誌情報はINSPIRE-HEPの現在の収録内容に従うため、原文の表記と異なる場合がある。
\end{itemize}
\end{document}
